%% LyX 2.3.6.1 created this file.  For more info, see http://www.lyx.org/.
%% Do not edit unless you really know what you are doing.
\documentclass[english]{article}
\usepackage[T1]{fontenc}
\usepackage[latin9]{inputenc}
\usepackage{babel}
\usepackage{array}
\usepackage{varioref}
\usepackage{multirow}
\usepackage{amsmath}
\usepackage{amsthm}
\usepackage{amssymb}

\makeatletter

%%%%%%%%%%%%%%%%%%%%%%%%%%%%%% LyX specific LaTeX commands.
%% Because html converters don't know tabularnewline
\providecommand{\tabularnewline}{\\}

%%%%%%%%%%%%%%%%%%%%%%%%%%%%%% Textclass specific LaTeX commands.
\numberwithin{equation}{section}
\numberwithin{figure}{section}
\usepackage[natbibapa]{apacite}
\theoremstyle{plain}
\newtheorem{thm}{\protect\theoremname}
\theoremstyle{definition}
\newtheorem{defn}[thm]{\protect\definitionname}
\theoremstyle{definition}
\newtheorem{example}[thm]{\protect\examplename}
\theoremstyle{plain}
\newtheorem{lem}[thm]{\protect\lemmaname}
\theoremstyle{plain}
\newtheorem{prop}[thm]{\protect\propositionname}

%%%%%%%%%%%%%%%%%%%%%%%%%%%%%% User specified LaTeX commands.
\usepackage{guessing}

\titletag{Accepted manuscript}
\makeatletter
\titlerunning{\@title}
\makeatother
\authorrunning{Jonas Moss}

\usepackage{booktabs}
%\renewcommand{\sqrt}[1]{{(#1)^{1/2}}}
\DeclareMathOperator{\tr}{tr}
\DeclareMathOperator{\Tr}{Tr}
\DeclareMathOperator{\Var}{Var}
\DeclareMathOperator{\sd}{sd}
\DeclareMathOperator{\argmin}{argmin}
\DeclareMathOperator{\Cor}{Cor}
\DeclareMathOperator{\Cov}{Cov}
\DeclareMathOperator{\diag}{diag} 
\DeclareMathOperator{\argmax}{argmax}
\DeclareMathOperator{\supp}{supp}

\usepackage{enumitem}
\setlist[enumerate]{label = (\roman*)}

\makeatother

\providecommand{\definitionname}{Definition}
\providecommand{\examplename}{Example}
\providecommand{\lemmaname}{Lemma}
\providecommand{\propositionname}{Proposition}
\providecommand{\theoremname}{Theorem}

\begin{document}
\title{Measuring agreement using guessing models and knowledge coefficients}
\author{Jonas Moss \orcid{0000-0002-6876-6964}\\
Department of Data Science and Analytics\\
BI Norwegian Business School\emph{}\\
\emph{jonas.moss@bi.no}}
\maketitle
\begin{abstract}
Several measures of agreement, such as the Perreault--Leigh coefficient,
the $\textsc{AC}_{1}$, and the recent coefficient of van Oest, are
based on explicit models of how judges make their ratings. To handle
such measures of agreement under a common umbrella, we propose a class
of models called guessing models, which contains most models of how
judges make their ratings. Every guessing model have an associated
measure of agreement we call the knowledge coefficient. Under certain
assumptions on the guessing models, the knowledge coefficient will
be equal to the multi-rater Cohen's kappa, Fleiss' kappa, the Brennan--Prediger
coefficient, or other less-established measures of agreement. We provide
several sample estimators of the knowledge coefficient, valid under
varying assumptions, and their asymptotic distributions. After a sensitivity
analysis and a simulation study of confidence intervals, we find that
the Brennan--Prediger coefficient typically outperforms the others,
with much better coverage under unfavorable circumstances. 
\end{abstract}

\section{Introduction\label{sec:Introduction}}

The most popular measures of agreement are chance-corrected. These
can usually be written on the form 
\begin{equation}
\frac{p_{a}-p_{ca}}{1-p_{ca}},\label{eq:chance-corrected measure of agrement}
\end{equation}
where $p_{a}$ is the percent agreement and $p_{ca}$ is a notion
of chance agreement. The best known coefficients in this class are
the (weighted) Cohen's kappa \citeyearpar{Cohen1960-mp,Cohen1968-sb},
Krippendorff's \citeyearpar{Krippendorff1970-gl} alpha, Scott's \citeyearpar{Scott1955-vn}
pi, and Fleiss' \citeyearpar{Fleiss1971-to} kappa. The difference
between these measures lies solely in their definition of the chance
agreement, $p_{ca}$. These coefficient make few to no assumptions
about the underlying distribution of ratings, and can be regarded
as non-parametric. 

It is also possible to model the judgment process directly, and then
attempt to derive reasonable chance-corrected measures of agreement
from these models \citep{Janes1979-ih}. Examples of measures of agreements
developed in this way include the Perreault--Leigh coefficient \citep{Perreault1989-mg},
the $\textsc{AC}_{1}$ \citep{Gwet2008-qk}, Maxwell's RE coefficient
\citep{Maxwell1977-lf}, Aickin's $\alpha$ \citep{Aickin1990-nj},
the estimators of \citet{Klauer1996-jd}, and the more recent coefficient
of van Oest \citep{Van_Oest2019-pm,Van_Oest2021-jt}. These measures
of agreement depend on the parameters of the underlying judgment process,
and may be considered semi-parametric instead of non-parametric. The
models used by the above-mentioned authors may be called \emph{guessing
models}, as they represent ratings as being either known or guessed. 

To make it clear what these models are about, consider the ``textbook
case argument'' of \citet{Grove1981-wb} \citep[see][Chapter 4, for an extended justification]{Gwet2014-iu}.
When two judges classify people into, say, psychiatric categories,
some people are bound to be ``textbook cases'', i.e., being classifiable
without much effort. Disagreement between competent judges will mostly
occur when subjects are hard to classify, when the judges have to
guess. But judges may agree on hard subjects as well, simply due to
chance. We can then define a coefficient of ``agreement due to knowledge''
as the proportion of textbook cases.

The guessing model, introduced in the next section, will encompass
the textbook case model and many more. As we will show, it is a generalization
of several judgment process models discussed in the literature on
measures of agreement. Any guessing model is associated with a \emph{knowledge
coefficient, }a measure of agreement defined directly from its parameters.
These coefficients generalize the ``agreement due to knowledge''
from Grove's textbook case to more general settings. The knowledge
coefficient can, under various additional assumptions, be easily estimated
from the data; the details are in Theorem \ref{thm:main theorem}.
In some cases, it equals already established coefficients such as
the Brennan--Prediger coefficient \citet{Brennan1981-kf} or Fleiss'
kappa, but we will establish some less familiar formulas as well.
We provide methods for doing inference for our proposed coefficients,
based on the delta rule and the theory of $U$-statistics. Using sensitivity
analyses and confidence interval simulations, we find that the Brennan--Prediger
coefficient generally outperforms its competitors as an estimator
of the knowledge coefficient, with reasonably small bias and variance
in a variety of circumstances. 

\section{Guessing models\label{sec:Guessing-models}}

We work in the setting where one rating is definitely true, such as
psychiatric diagnoses. Thus we exclude problems such as measuring
agreement between movie reviewers, where there is no true rating.
We also exclude measures of agreement between continuous measurement
instrument, as continuous ratings are rarely exactly right. For instance,
an instrument for measuring blood glucose may be decidedly better
than another, but will never be precisely on the spot.

We will consider only agreement studies with a rectangular design,
i.e., when $R$ judges rate $n$ items into one of $C<\infty$ categories,
with every item being rated exactly once by every judge. Moreover,
we will understand the set of judges as fixed and the set of items
as being random and increasing with $n$. These assumptions may not
be necessary for all of the results in this paper, but will make the
presentation easier to follow, and are necessary for the asymptotic
results. Denote the probability that the $R$ judges will rate an
item as belonging to the categories $x=(x_{1},x_{2},\ldots x_{R})$
by $p(x)$.

The joint distribution of the \emph{guessing model} is
\begin{equation}
p(x\mid s,x^{\star})=\prod_{r=1}^{R}\left[s_{r}1[x_{r}=x^{\star}]+(1-s_{r})q_{r}(x_{r})\right].\label{eq:full guessing model}
\end{equation}
where
\begin{itemize}
\item $x_{r}$ is the rating given by the $r$th judge on an item.
\item $s=\{s_{1},s_{2},\ldots,s_{r}\}$ are the \emph{skill-difficulty parameters},
the probabilities that the $r$th judge knows the correct classification
of an item. The skill-difficulty parameters can be deterministic or
random. For instance, they can be sampled from a Beta distribution.
\item $x^{\star}$ is the true classification the item, an unknown latent
variable. These are assumed to be independent of the skill-difficulty
parameters $s$. The distribution of $x^{\star}$ is $t(x^{\star})$,
the \emph{true rating distribution}.
\item $q_{r}(x)$ are the \emph{guessing distributions}, the distributions
the ratings are drawn from when the $r$th judge does not know the
true classification of the item.
\end{itemize}
We may use $t(x^{\star})$ to remove the dependence on $x_{i}^{\star}$
from the univariate guessing model, giving
\begin{equation}
p(x_{r}\mid s)=s_{r}t(x_{r})+(1-s_{r})q_{r}(x_{r}).\label{eq:univariate guessing model}
\end{equation}
The interpretation of $p(x_{r}\mid s)$ is straight-forward. When
faced with an item, a judge $r$ knows its true classification, drawn
from $t(x)$, with probability $s_{r}$. If the judge doesn't know
the true classification, the rating will be drawn at random from his
potentially idiosyncratic guessing distribution $q_{r}(x)$. We do
not allow for guessing distributions $q_{r}$ that depend both on
both the true classification $x^{\star}$ and the judge, as it would
make the parameters unidentifiable.

We have said nothing about the joint distribution of $(s_{1},s_{2},\dots,s_{R},x^{\star})$
except that $x^{\star}$ is independent of the skill-difficulty parameters.
This assumption is not realistic in all situations. For instance,
correctly diagnosing patients with Down syndrome is easier than correctly
diagnosing patients with ADHD, implying that $E[s_{r}\mid x^{\star}=\text{Down syndrome]}$
dominates $E[s_{r}\mid x^{\star}=\text{ADHD]}$, which violates independence
of $s_{r}$ and $x^{\star}$. The independence assumption is not needed
for the definition of the guessing model to make sense, but will be
used in the remainder of the paper as it is required for Theorem \ref{thm:main theorem}.

In most settings with latent parameters one would decide on a model
for them, such as multivariate normal in the case of linear random
effects models. Instead of following this route, we will impose additional
assumptions on the skill-difficulty parameters $s$, the guessing
distributions $q_{r}$, and/or the true distribution $t$ to make
the problem manageable. 

The guessing model \ref{eq:full guessing model} has, to our knowledge,
not been presented in this generality before. \citet[Theorem 5 and Section 9]{Klauer1996-jd}
define a model of almost as high generality, but does not allow the
the skill-difficulty parameters to differ between items.

\subsection{Knowledge coefficient}

We have introduced the guessing model in order to define the notion
of ``agreement due to knowledge'' in a precise way. To gain an intuition
about what we're getting at, first consider the case of two judges
with potentially different, but deterministic, skill-difficulty parameters
$s_{1}$ and $s_{2}$. The probability that two judges agree on the
classification of an item because they both know its classification
is the product of their skill parameters, or $\nu=s_{1}s_{2}$. As
``agree on classification of an item because they both know its classification''
is cumbersome to read, we will call it ``knowledgeable agreement''
or ``agree knowledgeably'' from now. Extending this notion to $R$
judges, $\nu=\binom{R}{2}^{-1}\sum_{r_{1}>r_{2}}s_{r_{1}}s_{r_{2}}$
is the probability that a randomly selected pair of judges will agree
knowledgeably on a pair of ratings. 

Another simple case happens when there are $R$ judges with random
skill-difficulty parameters that do not wary across judges when the
item is fixed, i.e., $S_{1}=S_{2}=\cdots=S_{R}$, where we use capital
letters to emphasize that the $S_{r}$ are random. Now $E(S_{r}^{2})$
is the probability of knowledgeable agreement. Finally, in the general
guessing model, we find that the probability of knowledgeable agreement
among two judges is
\begin{equation}
\nu=\binom{R}{2}^{-1}\sum_{r_{1}>r_{2}}E[S_{r_{1}}S_{r_{2}}].\label{eq:knowledgable agreement}
\end{equation}


\subsection{Earlier guessing models\label{subsec:Earlier-guessing-models}}

The guessing model and its associated knowledge coefficient are extensions,
formalizations, or slight modifications of models or coefficients
used in several earlier papers. 

\subsubsection{The two models of \citet{Maxwell1977-lf} \label{subsec:maxwell}}

\citet[Section 3]{Maxwell1977-lf} works in the setting of two judges
and binary ratings. From his Table II one can derive the the joint
model for two ratings $x_{1},x_{2}$ by two judges as

\begin{equation}
p(x_{1},x_{2})=\alpha p(x_{1})1[x_{1}=x_{2}]+(1-\alpha)p(x_{1})p(x_{2}),\label{eq:maxwell joint}
\end{equation}
where $p$ is the marginal distribution of the data. Maxwell then
shows that $\alpha=\Cor(X_{1},X_{2})$. 

Maxwell's joint distribution is the unconditional variant of a guessing
model (\ref{eq:full guessing model}) with two judges, i.e., a model
on the form 
\begin{align*}
p(x_{1},x_{2}\mid s,x^{\star}) & =\{s_{1}1[x_{1}=x^{\star}]+(1-s_{1})q_{r}(x_{1})\}\{s_{2}1[x_{2}=x^{\star}]+(1-s_{2})q_{r}(x_{2})\}
\end{align*}
with associated knowledge coefficient $\nu=\alpha$. The guessing
model satisfies 
\begin{enumerate}
\item The judges' guessing distributions are equal to the marginal distribution,
i.e., $q_{1}(x)=q_{2}(x)=p(x)$.
\item The true distribution is assumed to be equal to the marginal distribution,
i.e., $t(x)=p(x)$.
\item Both judges share the same skill-difficulty parameter $s$. It is
Bernoulli distributed with success probability $\alpha$, so that
$\alpha=P(s=1)=Es^{2}$. Then $s$ will $1$ if the the case is easy
to judge (i.e., a textbook case) and $0$ otherwise, and the probability
of an item being a textbook case is $\alpha$. 
\end{enumerate}
In Section 4, \citet{Maxwell1977-lf} is still working with binary
data and two judges. He describes a guessing model where (iii) above
still holds, but (i) and (ii) are replaced with
\begin{enumerate}
\item Both the judges' guessing distributions are uniform, i.e., $q_{1}(x)=q_{2}(x)=1/2$
and
\item The true distribution $t(x)$ is arbitrary.
\end{enumerate}
Then he derives the knowledge for this model, the Maxwell RE (abbreviation
of \emph{random error}) coefficient for binary data, a special case
of the Brennan--Prediger coefficient, 
\begin{equation}
\nu_{BP}=\frac{p_{a}-1/C}{1-1/C},\label{eq:Brennan-Prediger coefficient}
\end{equation}
where $C$, in this case equal to $2$, is the number of categories. 

\subsubsection{Perreault--Leigh coefficient (1989)}

\citet{Perreault1989-mg} devise an explicit model for the rating
procedure involving two judges and $C<\infty$ categories. Using an
index for reliability $s\in[0,1]$, they define the univariate model
\begin{equation}
p(x)=st(x)+(1-s)C^{-1}.\label{eq:Perreault--Leigh model}
\end{equation}
The model is similar to the second Maxwell model, except that the
skill-difficulty parameters are deterministic and constant across
judges and the number of categories is arbitrary. The guessing distributions
are uniform, $q_{1}=q_{2}=1/C$, and $t(x)$ is arbitrary. From this
model they derive that $s=\sqrt{(p_{a}-1/C)/(1-1/C)}=\sqrt{\nu_{BP}}$,
the square root of the Brennan--Prediger coefficient. Hence the knowledge
coefficient is $\nu=s^{2}.$ 

\subsubsection{Aickin's coefficient (1990)}

\citet{Aickin1990-nj} works in a setting of two judges. He defines
the joint model for two ratings $x_{1},x_{2}$ by two judges as
\begin{equation}
p(x_{1},x_{2})=(1-\alpha)q_{1}(x_{1})q_{2}(x_{2})+\alpha1[x_{1}=x_{2}]\frac{q_{1}(x_{1})q_{2}(x_{1})}{\sum_{x=1}^{C}q_{1}(x_{1})q_{2}(x_{1})},\label{eq:Aickin model}
\end{equation}
with the goal of doing inference on $\alpha$. He does this using
maximum likelihood, estimating the distributions $q_{1}$ and $q_{2}$
alongside $\alpha$.

Aickin's model is a guessing model and $\alpha=\nu$ is its knowledge
coefficient. The assumptions of the guessing model are:
\begin{enumerate}
\item As in the first Maxwell model, both judges share the same skill-difficulty
parameter $s$. It is Bernoulli distributed with success probability
$\alpha$, so that $\alpha=P(s=1)=Es^{2}$.
\item The judges' guessing distributions are arbitrary.
\item The true distribution is assumed to be equal to $t(x)=q_{1}(x)q_{2}(x)/\sum_{x=1}^{C}q_{1}(x)q_{2}(x)$.
\end{enumerate}
That Aicken's model is a guessing model satisfying conditions (i)--(iii)
is a direct consequence of the following fact. Whenever the number
of judges is two and the skill-difficulty parameter $s\sim\text{Bernoulli}(\alpha)$
is the same for both judges, the guessing model has unconditional
distribution
\begin{equation}
p(x_{1},x_{2})=\alpha1[x_{1}=x_{2}]t(x_{1})+(1-\alpha)q_{1}(x_{1})q_{2}(x_{2}).\label{eq:generalized aicken model-2}
\end{equation}
The details are in the appendix, p. \pageref{subsec:Proof-that-Aickin's}.

Assumption (iii) is not justified by Aickin, and does not appear to
be necessary. If we define the generalized Aicken model as the guessing
model satisfying only (i) and (ii) above, with an arbitrary number
of judges $R\geq2$, its parameters $(\nu,q_{1},q_{2},...q_{R})$
are identified when $C>2$. This can be shown following the arguments
laid out in the proof of Theorem 1 of \citet{Klauer1996-jd}.

\subsubsection{The Klauer--Batchelder model (1996)}

\citet{Klauer1996-jd} performs a detailed structural analysis of
the guessing model with identical skill-difficulty parameters. In
our notation, their equation 2 is
\begin{equation}
p(x\mid s,x^{\star})=s1[x=x^{\star}]+(1-s)q_{r}(x),\label{eq:klauer-batcheleder}
\end{equation}
where the guessing distributions $q_{1},q_{2}$ and the true distribution
$t$ are arbitrary. They show that, provided the number of judges
is equal to two, then \citep[eq. 3]{Klauer1996-jd} 
\begin{equation}
p(x_{1},x_{2})=p(x_{1})p(x_{2})+s^{2}t(x_{1})[1[x_{1}=x_{2}]-t(x_{2})],\label{eq:klauer-batchelder2}
\end{equation}
which does not depend directly on the guessing distributions $q_{r}(x)$.
Equation (\ref{eq:klauer-batchelder2}) provides a nice interpretation
of the skill-difficulty parameter $s$: The higher $s$ is, the more
weight will be on the main diagonal of the agreement matrix and less
on the off-diagonal elements. Moreover, in Theorem 5, they extend
equation \ref{eq:klauer-batchelder2} to the case of two judges with
different deterministic skill-difficulty parameters.

They show the model is identified when the $C>2$, and propose to
estimate it by maximum likelihood using the EM algorithm developed
by \citet{Hu1994-np}. In addition, they show that $s^{2}$ equals
Cohen's kappa when both guessing distributions are equal to the true
distribution; they also show that $s^{2}$ equals the Brennan--Prediger
coefficient when both distributions are uniform. We generalize these
results to arbitrary skill-difficulty parameters and an arbitrary
number of judges in Theorem \ref{thm:main theorem} below.

\subsubsection{van Oest's coefficient (2019)}

Modifying the setup of \citet{Perreault1989-mg}, \citet{Van_Oest2019-pm}
develops a guessing model for two judges and $C<\infty$ categories.
He assumes the guessing distributions are equal to the marginal distribution,
i.e., $q_{1}(x)=q_{2}(x)=p(x)$, and that the skill-difficulty coefficients
are deterministic and constant across judges. The marginal model becomes
\begin{equation}
p(x)=st(x)+(1-s)p(x).\label{eq:van Oest model}
\end{equation}
\citet{Van_Oest2019-pm} proceeds to show that $s$ equals the weighted
Scott's pi under these circumstances.

\section{The knowledge coefficient}

\subsection{Definitions}

Let $w(x,y)$ be an\emph{ agreement weighting function}. This is a
function of two arguments that satisfies $w(x,y)\leq1$ and equals
$1$ when $x=y$, i.e., $w(x,x)=1$. The purpose of this function
is to measure the degree of similarity between $x$ and $y$, where
$1$ is understood as the maximal degree of similarity. While there
are infinitely many weighting functions, only three are in widespread
use. The first is the \emph{nominal weight}, 
\[
w(x_{1},x_{2})=1[x_{1}=x_{2}]=\begin{cases}
1, & x_{1}=x_{2},\\
0, & \textrm{otherwise}.
\end{cases}
\]
With this function, similarity does not come in degrees, but it does
with the \emph{quadratic weighting function}, $w(x_{1},x_{2})=1-(x_{1}-x_{2})^{2}$.
The \emph{absolute value weighting function} (sometimes called the
\emph{linear} weighting function) measures the similarity between
$x,y$ using the absolute value, i.e., $w(x_{1},x_{2})=1-|x_{1}-x_{2}|$. 
\begin{defn}
Recall that $p_{r}(x)$ is the marginal distribution of ratings for
judge $r$, and let $x_{1}$ and $x_{2}$ be ratings by two different
judges $r_{1}$ and $r_{2}$. Define the \emph{weighted agreement}
as
\begin{eqnarray}
p_{wa} & = & \binom{R}{2}^{-1}\sum_{r_{1}>r_{2}}\sum_{x_{1},x_{2}}w(x_{1},x_{2})p(x_{r_{1}},x_{r_{2}}),\label{eq:Weighted probability of agreement}
\end{eqnarray}
the \emph{weighted Cohen-type chance agreement} as
\begin{eqnarray}
p_{wc} & = & \binom{R}{2}^{-1}\sum_{r_{1}>r_{2}}\sum_{x_{1},x_{2}}w(x_{1},x_{2})p(x_{r_{1}})p(x_{r_{2}}),\label{eq:Fleiss-type weighted probability of chance agreement}
\end{eqnarray}
and the \emph{weighted Fleiss-type chance agreement }as
\begin{eqnarray}
p_{wf} & = & R^{-2}\sum_{r_{1},r_{2}}\sum_{x_{1},x_{2}}w(x_{1},x_{2})p(x_{r_{1}})p(x_{r_{2}}).\label{eq:Cohen-type weighted probability of chance agreement}
\end{eqnarray}
\end{defn}

The difference between the two notions of weighted chance agreement
should be clear enough. The Fleiss-type probability of chance agreement
counts the cases when a judge agrees with himself, while the Cohen-type
does not. 

Letting $\{x_{1},x_{2},...,x_{C}\}$ be the set of possible ratings
and $w$ an agreement weighting function, define the weighting matrix
$W$ as the $C\times C$ matrix with elements $W_{ir}=w(x_{i},x_{r})$.
Using $W$, we can write $p_{wf}=p^{T}Wp,$ where $p=R^{-1}\sum_{r}p_{r}$
is the marginal distribution of ratings, and $p_{wc}=\binom{R}{2}^{-1}\sum_{r_{1}>r_{2}}p_{r_{1}}^{T}Wp_{r_{2}}$.

When $w$ is the nominal weight, $p_{wa},p_{wc}$, and $p_{wf}$ are
proper probabilities, and are often referred to as unweighted probabilities
of (chance) agreement. We do not require that the weighting functions
to be non-negative, hence the quantities $p_{wa},p_{wc}$, and $p_{wf}$
are not, in general, proper probabilities. Since the number of categories
$C$ is finite, however, we may assume that the weighting function
is positive by normalizing, i.e., redefining the weighting function
to $1-[1-w(x_{1},x_{2})]/\max_{x_{1},x_{2}}(1-w(x_{1},x_{2}))$. 

\subsection{The knowledge coefficient theorem}

As we have seen, well-known coefficients such as Scott's pi and the
Brennan--Prediger can be understood as knowledge coefficients, albeit
under restrictive assumptions such as all judges being equally skilled.
The following theorem describes less stringent sets of assumptions
that keeps interpretable expressions for the knowledge coefficient.
We must assume either that all guessing distributions are equal or
that the skill-difficulty parameters have zero pairwise covariance
to get anywhere. In addition, we must assume something about either
the true distribution or the guessing distribution. We never have
to assume that every judge is equally competent, and we never have
to assume anything about the number of categories rated, however.
The content of the theorem, including the required assumptions, is
summarized in Table \ref{tab:coefficients}.

\begin{table}
\noindent \begin{centering}
\caption{\label{tab:coefficients}Coefficients covered in this paper}
\begin{tabular}{ccc}
\toprule\textbf{Assumptions} & \textbf{Name} & \textbf{Definition}\tabularnewline
 &  & \tabularnewline
\multicolumn{3}{c}{\textbf{Provided that all guessing distributions are equal}}\tabularnewline
 &  & \tabularnewline
\cmidrule[0.5pt]{2-3}Guessing $=$ Uniform & Brennan--Prediger coefficient & $\nu_{BP}=\frac{p_{a}-1/C}{1-1/C}$\tabularnewline
 &  & \tabularnewline
Guessing $=$ Marginal$^{*}$ or & Weighted Fleiss' kappa & $\nu_{F}=\frac{p_{wa}-p_{wf}}{1-p_{wf}}$\tabularnewline
\multirow{2}{*}{Guessing $=$ True$^{*}$ or} &  & \tabularnewline
 & Weighted Cohen's kappa & $\nu_{C}=\frac{p_{wa}-p_{wc}}{1-p_{wc}}$\tabularnewline
\multirow{2}{*}{True $=$ Marginal$^{*}$} &  & \tabularnewline
 & \emph{Cohen--Fleiss coefficient} & $\nu_{CF}=\frac{p_{wa}-p_{wc}}{1-p_{wf}}$\tabularnewline
 &  & \tabularnewline
\multicolumn{3}{c}{\textbf{Provided $E[s_{r_{1}}s_{r_{2}}]=E[s_{r_{1}}]E[s_{r_{2}}]$
for all $r_{1},r_{2}$}}\tabularnewline
 &  & \tabularnewline
\cmidrule[0.5pt]{2-3}True $=$ Marginal & \multirow{1}{*}{\emph{Cohen--Fleiss coefficient}} & \multirow{1}{*}{$\nu_{CF}=\frac{p_{wa}-p_{wc}}{1-p_{wf}}$}\tabularnewline
 &  & \tabularnewline
True $=$ Uniform & \multirow{1}{*}{\emph{Cohen--Brennan--Prediger coefficient}} & \multirow{1}{*}{$\nu_{CBP}=\frac{p_{wa}-p_{wc}}{1-\boldsymbol{1}^{T}W\boldsymbol{1}/C^{2}}$}\tabularnewline
 &  & \tabularnewline
\bottomrule &  & \tabularnewline
\end{tabular}
\par\end{centering}
\vspace{1mm}

\emph{\footnotesize{}Note}{\footnotesize{}: New coefficients in italics.}{\footnotesize\par}

{\footnotesize{}$^{*}$These three conditions are equivalent, see
(i) of Theorem \ref{thm:main theorem}.}{\footnotesize\par}
\end{table}

\begin{thm}[Knowledge Coefficient Theorem]
\label{thm:main theorem}Let $w$ be any agreement weighting function
and $W$ its associated agreement weighting matrix. Then the following
holds:
\begin{enumerate}
\item Assume all guessing distributions are equal, i.e, $q_{r}(x)=q(x)$.
Then the following are equivalent:
\[
q(x)=t(x),\quad p(x)=t(x),\quad q(x)=p(x)
\]
Assuming either of these,$p_{wc}=p_{wf}$, and the knowledge coefficient
equals both the weighted multi-rater Cohen's kappa (Conger's kappa)
and the weighted Fleiss' kappa
\[
\nu=\nu_{C}=\frac{p_{wa}-p_{wc}}{1-p_{wc}};\quad\nu=\nu_{F}=\frac{p_{wa}-p_{wf}}{1-p_{wf}}.
\]
\item Assume all guessing distributions $q_{r}$ are uniform. Then the knowledge
coefficient equals the Brennan--Prediger coefficient,
\[
\nu=\nu_{BP}=\frac{p_{a}-1/C}{1-1/C},
\]
where $p_{a}$ denotes the weighted agreement with nominal weights.
\item Assume that the skill-difficulty parameters have pairwise covariance
equal to zero, i.e., $E[s_{r_{1}}s_{r_{2}}]=E[s_{r_{1}}]E[s_{r_{2}}]$
for all $r_{1},r_{2}$. Then the knowledge coefficient equals
\[
\nu=\frac{p_{wa}-p{}_{wc}}{1-t^{T}Wt}.
\]
In particular, if $t(x)=p(x)$, it equals the ``Cohen--Fleiss''
coefficient
\[
\nu=\nu_{CF}=\frac{p_{wa}-p_{wc}}{1-p_{wf}}.
\]
Moreover, if $t$ is uniform, it equals the ``Cohen--Brennan--Prediger''
coefficient,
\[
\nu=\nu_{CBP}=\frac{p_{wa}-p_{wc}}{1-\boldsymbol{1}^{T}W\boldsymbol{1}/C^{2}}.
\]
\end{enumerate}
\end{thm}

\begin{proof}
The knowledge coefficient theorem is proved in the appendix, page
\pageref{subsec:proof of knowledge theorem}.
\end{proof}
Every coefficient in the theorem can be estimated by substituting
the values of $p_{wa}$, $p_{wf}$, and $p_{wc}$ for their sample
variants. Under any of the equivalent conditions of Theorem \ref{thm:main theorem}
part (i), we have that that $\nu$ equals 
\[
\nu_{F}=\frac{p_{wa}-p_{wf}}{1-p_{wf}},\quad\nu_{C}=\frac{p_{wa}-p_{wc}}{1-p_{wc}},\quad\nu_{CF}=\frac{p_{wa}-p_{wc}}{1-p_{wf}}.
\]
The coefficient $\nu_{F}=\frac{p_{wa}-p_{wf}}{1-p_{wf}}$ is a weighted
Fleiss' kappa. This coefficient is strongly related to the weighted
Krippendorff's alpha, which we denote by $\hat{\alpha}$. Indeed,
it is easy to see that $\hat{\alpha}$ is a linear transformation
of $\hat{\nu}_{F}$,
\begin{equation}
\hat{\alpha}=\hat{\nu}_{F}+\frac{1}{N}(1-\hat{\nu}_{F}),\label{eq:krippendorffs alpha}
\end{equation}
where $N$ is the total number of ratings made \citep[see the appendix of][]{Anon2022-ep}
and $\hat{\nu}_{F}$, the sample weighted Fleiss kappa. Thus $\hat{\alpha}$
is a consistent estimator of $\nu_{F}$. 

On the other hand, $\nu_{C}$ is a weighted multi-rater Cohen's kappa
of the form discussed by \citet{Berry1988-nc} and \citet{Janson2001-pk};
it can also be regarded as a weighted Conger's kappa \citep{Conger1980-yx}.
I have not seen anything like $\nu_{CF}$, a curious combination of
Cohen's kappa and Fleiss' kappa, probably because it is not a classical
chance-corrected chance measure of agreement, as its numerator chance
agreement $p_{wc}$ is distinct from the denominator chance agreement
$p_{wf}$. The required condition for the Cohen--Fleiss coefficient,
$p(x)=t(x)$, holds if and only if
\[
R^{-1}\sum_{r=1}^{R}\left[(1-s_{r})q_{r}(x)+s_{r}t(x)\right]=t(x)\Longleftrightarrow\sum_{r=1}^{R}(1-s_{r})q_{r}(x)=t(x)\sum_{r=1}^{R}(1-s_{r}),\quad\text{for all }r.
\]
which hols trivially if $q_{r}(x)=t(x)$ for all $r$. Hence the Cohen--Fleiss
coefficient is consistent for the population knowledge coefficient
under strictly more situations than weighted Fleiss' kappa and weighted
multi-rater Cohen's kappa, and may be preferred to them by a researcher
who buys the rationale behind the guessing model. 

The Cohen-Brennan--Prediger coefficient will be equal to the knowledge
coefficient if the true distribution equals the uniform distribution
and the skill-difficulty parameters have pairwise covariances equal
to zero. This scenario may be uncommon, but can happen if the designer
of the agreement study has complete control over the true ratings. 

Part (ii) of Theorem \ref{thm:main theorem} concerns the case when
a judge who does not know the correct classification of an item guesses
uniformly at random, so that $q_{r}(x)=C^{-1}$ for all judges $r$.
This assumption is used by e.g. \citet{Brennan1981-kf}, \citet{Maxwell1977-lf},
\citet{Gwet2008-qk} and \citet{Perreault1989-mg}. 
\begin{example}
\label{exa:Zapf exampe} \citet{Zapf2016-eu} did a case study on
histopathological assessment of breast cancer. The number of judges
was $R=4$, the number breast cancer biopsies rated was $n=50,$and
the number of categories $C=5$. The estimated coefficients are
\[
\hat{\nu}_{CF}=0.574,\quad\hat{\nu}_{F}=0.562,\quad\hat{\nu}_{C}=0.567,\quad\hat{\nu}_{BP}=0.604,\quad\hat{\nu}_{CBP}=0.519.
\]
In this case, the coefficients are quite close, suggesting that the
guessing distributions are close to the marginal distribution and
that the marginal distribution is close to the uniform distribution.
The observed marginal distribution in this data set $0.255,0.025,0.12,0.21,0.39$.
This appears to be quite far away from the uniform distribution, which
raises the question of how sensitive the Brennan--Prediger coefficient
is to the uniformity assumption.
\end{example}


\section{Sensitivity and performance}

Recall the assumptions on the coefficients of Theorem \ref{thm:main theorem}
\begin{itemize}
\item \textbf{Brennan--Prediger}: All guessing distributions are equal
to the uniform distribution.
\item \textbf{Cohen's kappa, Fleiss' kappa}: All guessing distributions
are equal, the true distribution equals the marginal distribution.
\item \textbf{Cohen--Fleiss}: The true distribution equals the marginal
guessing distribution, $E[s_{r_{1}}s_{r_{2}}]=E[s_{r_{1}}]E[s_{r_{2}}]$
for all $r_{1},r_{2}$.
\item \textbf{Cohen--Brennan--Prediger}: The true distribution is uniform,
$E[s_{r_{1}}s_{r_{2}}]=E[s_{r_{1}}]E[s_{r_{2}}]$ for all $r_{1},r_{2}$.
\end{itemize}
These assumptions are quite stringent, and will realistically never
hold exactly. In this section, we do two sensitivity--performance
studies to check how well the coefficients perform when the assumptions
are broken. Theorem \ref{thm:main theorem} contains two classes of
coefficients. The first class, containing the Brennan--Prediger coefficient
in addition to Cohen's and Fleiss' kappa, requires at minimum that
all guessing distributions are equal. The second class, containing
the Cohen--Fleiss and Cohen--Brennan--Prediger coefficient, requires
that all pairs of skill-difficulty parameters have zero covariance.
We will do two studies, one where $s_{r_{1}}$ and $s_{r_{2}}$ have
zero covariance and one where $s_{r_{1}}$ and $s_{r_{2}}$ are correlated.
We restrict our study the the case of nominal weights. 

\subsection{When $E[s_{r_{1}}s_{r_{2}}]=E[s_{r_{1}}]E[s_{r_{2}}]$ \label{subsec:sensitivity independence}}

We will use the a special case of the guessing model we call the \emph{judge
skill model}. In this model the skill-difficulty parameter $s$ is
deterministic. This is a generalization of the models used by e.g.
\citet{Perreault1989-mg} and \citet{Van_Oest2019-pm} to allow for
judges with different skill levels. Since $s$ is deterministic, $E[s_{r_{1}}s_{r_{2}}]=E[s_{r_{1}}]E[s_{r_{2}}]$
and the main condition of Theorem \ref{thm:main theorem} part (iii)
is satisfied. Under the judge skill model, it is fairly easy to calculate
the theoretical values of the five coefficients in Theorem \ref{thm:main theorem}.
To calculate $p_{wa}$, we use representation (i) of Lemma \ref{lem:expression lemma}
(p. \pageref{lem:expression lemma}), in the appendix. Since $p_{wf}=p^{T}Wp$,
where $p$ is the marginal distribution, and $p_{wc}=\binom{R}{2}^{-1}\sum_{r_{1}>r_{2}}p_{r_{1}}^{T}Wp_{r_{2}}$,
the values of $p_{wf}$ and $p_{wc}$ are easily calculated.

The parameters $R$, $C$ and $s$ are sampled as follows:
\begin{enumerate}
\item The number of judges ($R$) is sampled uniformly from $[2,20]$. 
\item The number of categories ($C$) is sampled uniformly from $[2,10]$.
\item The $R$ skill--difficulty parameters $s_{1},...,s_{R}$ are drawn
independently from a beta distribution with parameters $7$ and $1.5$.
This is a slightly dispersed, asymmetric distribution with a mean
of $0.82$. 
\end{enumerate}
We study what happens when the true distribution deviates from the
uniform (an assumption for $\nu_{CBP}$) and / or the guessing distribution
deviates from the uniform distribution (an assumption for $\nu_{BP}$).
The numbers are the simulated mean absolute deviations from the true
knowledge coefficient $E|\nu-\nu_{x}|$, where $x$ is one of $F,C,BP,CF,$or
$CBP$. The smallest numbers by orders of magnitude on each row is
in bold.

\subsubsection{True distribution centered on the uniform distribution\label{subsec:True-distribution-centered-uniform}}

If the variability of the true distribution is ``None'', it equals
the uniform distribution. If the variability is ``Low'', it is sampled
from a symmetric Dirichlet distribution \citep[Chapter 49]{Johnson1994-ag}
with concentration parameter $\alpha=10$; if variability is ``High'',
$\alpha=0.5$. Likewise, for the guessing distributions, if the variability
is ``None'', all guessing distributions are equal to the uniform
distribution. If the variability is ``Low'', the guessing distributions
are sampled from a symmetric Dirichlet distribution with concentration
parameter $\alpha=10$. Finally, if the variability is ``High'',
they are sampled from a symmetric Dirichlet distribution with $\alpha=0.5$.

\begin{table}
\caption{\label{tab:centered on uniform}Sensitivity analysis when $E[s_{r_{1}}s_{r_{2}}]=E[s_{r_{1}}]E[s_{r_{2}}]$.
True distribution centered on the uniform distribution.}

\noindent \begin{centering}
\begin{tabular}{ccccccccc}
\toprule & \multicolumn{2}{c}{Variability} &  & \multicolumn{5}{c}{Coefficient}\tabularnewline
\cmidrule[0.5pt]{2-3}\cmidrule[0.5pt]{5-9} & Guessing$^{\star}$ & True$^{*}$ &  & Cohen-Fleiss & Fleiss & Cohen & BP & Cohen--BP\tabularnewline
\midrule & \multirow{3}{*}{None} & None &  & \textbf{0.0e+00} & \textbf{0.0e+00} & \textbf{0.0e+00} & \textbf{0.0e+00} & \textbf{0.0e+00}\tabularnewline
 &  & Low &  & 3.8e-03 & 3.8e-03 & 3.8e-03 & \textbf{0.0e+00} & 1.2e-02\tabularnewline
 &  & High &  & 7.5e-02 & 7.5e-02 & 7.5e-02 & \textbf{0.0e+00} & 1.7e-01\tabularnewline
 & \multirow{3}{*}{Low} & None &  & \textbf{5.2e-05} & \textbf{4.9e-05} & \textbf{4.3e-05} & \textbf{5.5e-05} & \textbf{1.7e-05}\tabularnewline
 &  & Low &  & 3.8e-03 & 3.8e-03 & 3.8e-03 & \textbf{5.7e-04} & 1.2e-02\tabularnewline
 &  & High &  & 7.4e-02 & 7.4e-02 & 7.4e-02 & \textbf{2.0e-03} & 1.7e-01\tabularnewline
 & \multirow{3}{*}{High} & None &  & \textbf{7.3e-04} & \textbf{8.0e-04} & \textbf{6.9e-04} & \textbf{8.7e-04} & \textbf{2.7e-04}\tabularnewline
 &  & Low &  & \textbf{3.5e-03} & \textbf{4.2e-03} & \textbf{3.9e-03} & \textbf{2.4e-03} & 1.2e-02\tabularnewline
 &  & High &  & 7.3e-02 & 7.4e-02 & 7.4e-02 & \textbf{7.9e-03} & 1.7e-01\tabularnewline
\bottomrule &  &  &  &  &  &  &  & \tabularnewline
\end{tabular}
\par\end{centering}
\vspace{1mm}

{\scriptsize{}$^{*}$Variability of the true distributions: Baseline:
True distribution is uniform.}{\scriptsize\par}

{\scriptsize{}$^{\star}$Variability of the guessing distributions.
Baseline: All guessing distributions are equal to the true distribution. }{\scriptsize\par}
\end{table}

From Table \ref{tab:centered on uniform} we see that the Cohen--Brennan--Prediger
coefficient performs worst in every setting, , usually by a large
margin, except when the true distribution is uniform. Moreover, the
Brennan--Prediger coefficient performs well in every scenario. The
bias $|\nu_{BP}-\nu|$ will be likely be overshadowed by sampling
variability for most conceivable sample sizes. Finally, there is little
difference between Cohen's kappa, Fleiss' kappa, and the Cohen--Fleiss
coefficient. Their biases are quite small, at least when the true
distribution isn't far away from the uniform distribution.

\subsubsection{True distribution centered on the marginal distribution\label{subsec:True-distribution-centered-marginal}}

To derive Fleiss' kappa, Cohen's kappa, and the Cohen--Fleiss coefficient,
we assumed that the true distribution equals the marginal distribution.
We can use an asymmetric Dirichlet distribution to extend this scenario,
just as we used the symmetric Dirichlet distribution to extend the
scenario when the true distribution is uniform. This time we won't
use the uniform distribution as a base true distribution, but randomly
generated distribution $h$, from a symmetric Dirichlet distribution
with $\alpha=5$, instead. The rest of the settings are identical
to the previous sensitivity study.

\begin{table}
\caption{\label{tab:centered on marginal}Sensitivity analysis when $E[s_{r_{1}}s_{r_{2}}]=E[s_{r_{1}}]E[s_{r_{2}}]$.
True distribution centered on the marginal guessing distribution.}

\noindent \begin{centering}
\begin{tabular}{ccccccccc}
\toprule & \multicolumn{2}{c}{Variability} &  & \multicolumn{5}{c}{Coefficient}\tabularnewline
\cmidrule[0.5pt]{2-3}\cmidrule[0.5pt]{5-9} & Guessing$^{\star}$ & True$^{*}$ &  & Cohen-Fleiss & Fleiss & Cohen & BP & Cohen--BP\tabularnewline
\midrule & \multirow{3}{*}{None} & None &  & \textbf{0.0e+00} & \textbf{0.0e+00} & \textbf{0.0e+00} & 1.1e-02 & 2.5e-02\tabularnewline
 &  & Low &  & \textbf{2.1e-02} & \textbf{2.1e-02} & \textbf{2.1e-02} & \textbf{1.1e-02} & \textbf{8.2e-02}\tabularnewline
 &  & High &  & 3.4e-01 & 3.4e-01 & 3.4e-01 & \textbf{1.8e-02} & 4.7e-01\tabularnewline
 & \multirow{3}{*}{Low} & None &  & \textbf{1.4e-04} & \textbf{4.6e-04} & \textbf{3.6e-04} & 1.3e-02 & 3.0e-02\tabularnewline
 &  & Low &  & \textbf{2.2e-02} & \textbf{2.2e-02} & \textbf{2.2e-02} & \textbf{1.4e-02} & \textbf{9.1e-02}\tabularnewline
 &  & High &  & 3.2e-01 & 3.3e-01 & 3.3e-01 & \textbf{1.8e-02} & 4.6e-01\tabularnewline
 & \multirow{3}{*}{High} & None &  & \textbf{1.1e-03} & \textbf{3.5e-03} & \textbf{2.8e-03} & 3.0e-02 & 7.4e-02\tabularnewline
 &  & Low &  & \textbf{2.5e-02} & \textbf{2.7e-02} & \textbf{2.6e-02} & \textbf{3.0e-02} & 1.3e-01\tabularnewline
 &  & High &  & 3.2e-01 & 3.2e-01 & 3.2e-01 & \textbf{3.4e-02} & 4.7e-01\tabularnewline
\bottomrule &  &  &  &  &  &  &  & \tabularnewline
\end{tabular}
\par\end{centering}
\vspace{1mm}

{\scriptsize{}$^{*}$Variability of the true distributions: Baseline:
True distribution equals the marginal guessing distribution.}{\scriptsize\par}

{\scriptsize{}$^{\star}$Variability of the guessing distributions.
Baseline: All guessing distributions are equal to the marginal guessing
distribution.}{\scriptsize\par}
\end{table}

The results are in Table \ref{tab:centered on marginal}. We see that
the Cohen--Brennan--Prediger coefficient performs worst in every
setting, usually by a large margin. Surprisingly, the Brennan--Prediger
coefficient performs best in $6/9$ cases, and its performance is
good in the remaining cases too. Some of the biases in the table are
unacceptably large. Only the Brennan--Prediger coefficient has a
bias less than $0.1$ in every case. Finally, there is little difference
between Cohen's kappa, Fleiss' kappa, and the Cohen--Fleiss coefficient.
The Cohen--Fleiss coefficient does better, especially when the true
distribution equals the marginal distribution, but the difference
in performance is insignificant under slight deviations from equality.

\subsection{When $E[s_{r_{1}}s_{r_{2}}]\protect\neq E[s_{r_{1}}]E[s_{r_{2}}]$\label{subsec:dependence}}

To model dependent skill-difficulty parameters, let $F$ be a $R$-variate
distribution with $s\sim F$. Evidently, the only restriction on $F$
is that it's a multivariate distribution function on $[0,1]$. A natural
way to model situation is to use copulas for the dependence structure
and a density on $[0,1]$ for the marginals \citep{Nelsen2007-qj}.
We will use the $R$-variate Gaussian copula with uniform correlation
structure, i.e., the correlation matrix parameter
\[
\left[\begin{array}{cccc}
1 & \rho & \cdots & \rho\\
\rho & 1 & \cdots & \rho\\
\vdots & \vdots & \ddots & \vdots\\
\rho & \rho & \cdots & 1
\end{array}\right].
\]
Denote this Gaussian copula by $C_{\rho}$. Let $H_{(a,b)}$ be the
cumulative distribution function of a beta distribution with parameters
$a$ and $b$, and define 
\[
F_{\rho,a,b}(s_{1},s_{2},\ldots,s_{R})=C_{\rho}(H_{(a,b)}(s_{1}),H_{(a,b)}(s_{2}),\ldots,H_{(a,b)}(s_{R})).
\]
Then $F_{\rho,a,b}$ is a reasonable model for dependent skill-difficulty
parameters.

We use the small correlation parameter $\rho=0.2$. The rest of the
settings are exactly the same as the previous setting, including the
beta parameters $a=7,b=1/2$. We see that the Cohen--Brennan--Prediger
coefficient still outperforms the alternatives when the true distribution
equals the uniform distribution, but only by a an order of magnitude.
Since the alternatives to the Cohen--Brennan--Prediger coefficient
also performs well in this situation, with biases likely to be overshadowed
by sampling error, the case for it remains weak. Similarly, we see
that the Cohen--Fleiss coefficient still outperforms Cohen's kappa
and Fleiss' kappa, roughly by one order of magnitude, a much smaller
margin than before. These comments also hold for the case of $\rho=0.5,0.9$,
which can be found in the online appendix (p. \pageref{tab:centered on uniform-1-1}).

\begin{table}
\caption{\label{tab:centered on uniform0.2}Sensitivity analysis when $E[s_{r_{1}}s_{r_{2}}]\protect\neq E[s_{r_{1}}]E[s_{r_{2}}]$.
True distribution centered on the uniform distribution.}

\noindent \begin{centering}
\begin{tabular}{ccccccccc}
\toprule & \multicolumn{2}{c}{Variability} &  & \multicolumn{5}{c}{Coefficient}\tabularnewline
\cmidrule[0.5pt]{2-3}\cmidrule[0.5pt]{5-9} & Guessing$^{\star}$ & True$^{*}$ &  & Cohen--Fleiss & Fleiss & Cohen & BP & Cohen--BP\tabularnewline
\midrule & \multirow{3}{*}{None} & None &  & \textbf{0.0e+00} & \textbf{0.0e+00} & \textbf{0.0e+00} & \textbf{0.0e+00} & 0.0e+00\tabularnewline
 &  & Low &  & 3.8e-03 & 3.8e-03 & 3.8e-03 & \textbf{0.0e+00} & 1.2e-02\tabularnewline
 &  & High &  & 7.4e-02 & 7.4e-02 & 7.4e-02 & \textbf{0.0e+00} & 1.7e-01\tabularnewline
 & \multirow{3}{*}{Low} & None &  & 4.6e-05 & 4.0e-05 & 3.4e-05 & 4.5e-05 & \textbf{8.7e-06}\tabularnewline
 &  & Low &  & 3.5e-03 & 3.6e-03 & 3.6e-03 & \textbf{5.6e-04} & 1.1e-02\tabularnewline
 &  & High &  & 7.8e-02 & 7.8e-02 & 7.8e-02 & 2.4e-03 & 1.7e-01\tabularnewline
 & \multirow{3}{*}{High} & None &  & 6.3e-04 & 5.8e-04 & 4.6e-04 & 6.2e-04 & \textbf{1.2e-04}\tabularnewline
 &  & Low &  & 3.6e-03 & 4.3e-03 & 4.0e-03 & \textbf{2.2e-03} & 1.1e-02\tabularnewline
 &  & High &  & 8.2e-02 & 8.4e-02 & 8.3e-02 & \textbf{8.6e-03} & 1.8e-01\tabularnewline
\bottomrule &  &  &  &  &  &  &  & \tabularnewline
\end{tabular}
\par\end{centering}
\vspace{1mm}

{\scriptsize{}$^{*}$Variability of the true distributions: Baseline:
True distribution is uniform.}{\scriptsize\par}

{\scriptsize{}$^{\star}$Variability of the guessing distributions.
Baseline: All guessing distributions are equal to the true distribution.}{\scriptsize\par}
\end{table}

\begin{table}
\caption{\label{tab:centered on marginal0.2}Sensitivity analysis when $E[s_{r_{1}}s_{r_{2}}]\protect\neq E[s_{r_{1}}]E[s_{r_{2}}]$.
True distribution centered on the marginal guessing distribution.}

\noindent \begin{centering}
\begin{tabular}{ccccccccc}
\toprule & \multicolumn{2}{c}{Variability} &  & \multicolumn{5}{c}{Coefficient}\tabularnewline
\cmidrule[0.5pt]{2-3}\cmidrule[0.5pt]{5-9} & Guessing$^{\star}$ & True$^{*}$ &  & Cohen--Fleiss & Fleiss & Cohen & BP & Cohen--BP\tabularnewline
\midrule & \multirow{3}{*}{None} & None &  & \textbf{0.0e+00} & \textbf{0.0e+00} & \textbf{0.0e+00} & 1.0e-02 & 2.3e-02\tabularnewline
 &  & Low &  & 2.2e-02 & 2.2e-02 & 2.2e-02 & \textbf{1.2e-02} & 8.4e-02\tabularnewline
 &  & High &  & 3.2e-01 & 3.2e-01 & 3.2e-01 & \textbf{1.8e-02} & 4.5e-01\tabularnewline
 & \multirow{3}{*}{Low} & None &  & \textbf{7.5e-05} & 3.9e-04 & 2.9e-04 & 1.3e-02 & 3.0e-02\tabularnewline
 &  & Low &  & \textbf{2.3e-02} & \textbf{2.3e-02} & \textbf{2.3e-02} & 1.4e-02 & 9.1e-02\tabularnewline
 &  & High &  & 3.3e-01 & 3.3e-01 & 3.3e-01 & \textbf{2.1e-02} & 4.6e-01\tabularnewline
 & \multirow{3}{*}{High} & None &  & \textbf{6.2e-04} & 3.3e-03 & 2.4e-03 & 3.0e-02 & 7.1e-02\tabularnewline
 &  & Low &  & \textbf{2.4e-02} & 2.7e-02 & 2.6e-02 & 3.1e-02 & 1.3e-01\tabularnewline
 &  & High &  & 3.5e-01 & 3.5e-01 & 3.5e-01 & \textbf{3.6e-02} & 4.8e-01\tabularnewline
\bottomrule &  &  &  &  &  &  &  & \tabularnewline
\end{tabular}
\par\end{centering}
\vspace{1mm}

{\scriptsize{}$^{*}$Variability of the true distributions: Baseline:
True distribution equals the marginal guessing distribution.}{\scriptsize\par}

{\scriptsize{}$^{\star}$Variability of the guessing distributions.
Baseline: All guessing distributions are equal.}{\scriptsize\par}
\end{table}


\subsection{Recommendations}

We can draw three tentative conclusions from the sensitivity analysis.
\begin{enumerate}
\item Unless the researcher can make sure the true distribution is exactly
uniform, the Cohen--Brennan--Prediger coefficient is probably not
worth reporting. Its bias is often unacceptably large, frequently
larger than $0.1$.
\item The Brennan--Prediger coefficient performs reasonably well in all
situations, with biases less than $0.01$ when the true distribution
is centered on the uniform distribution. It performs decently in the
second scenario as well, with biases at around $0.05$ across the
board.
\item The Cohen--Fleiss coefficient does slightly better than Cohen's kappa
and Fleiss' kappa, but not enough to be important. 
\end{enumerate}
Since the Cohen--Fleiss coefficient does slightly better than Fleiss'
kappa and Cohen's kappa, it appears prudent to report it. However,
we believe it would be best not to. For both Fleiss' kappa and Cohen's
kappa are well-known chance-corrected measures of agreement. In contrast,
the Cohen--Fleiss coefficient is neither well-known nor a chance-corrected
measure of agreement. We recommend that you report Cohen's kappa (or
Fleiss' kappa) together with the Brennan--Prediger coefficient. This
solution accounts both for the scenario when the marginal distribution
is close to the true distribution and the scenario when the uniform
distribution is close to the marginal distribution.

\section{Inference}

The asymptotic distribution of the coefficients Table \ref{tab:coefficients}
can readily be calculated using the theory of $U$-statistics. 

The following Lemma is instrumental in constructing the confidence
intervals.
\begin{lem}
Define the parameter vectors $\boldsymbol{p}=(p_{wa},p_{wc},p_{wf})$
and $\hat{\boldsymbol{p}}=(\hat{p}_{wa},\hat{p}_{wc},\hat{p}_{wf})$,
and let $\Sigma$ be the covariance matrix with elements 
\begin{alignat*}{6}
\sigma_{11} & = & \sigma_{A}^{2} & = & \Var\mu_{wa}(\boldsymbol{X}_{1}) & ,\quad & \sigma_{12} & = & \sigma_{AC}^{2} & = & 2\Cov\left(\mu_{wa}(\boldsymbol{X}_{1}),\mu_{wc}(\boldsymbol{X}_{1})\right),\\
\sigma_{22} & = & \sigma_{C}^{2} & = & 4\Var\mu_{wc}(\boldsymbol{X}_{1}) & ,\quad & \sigma_{13} & = & \sigma_{AF}^{2} & = & 2\Cov\left(\mu_{wa}(\boldsymbol{X}_{1}),\mu_{wc}(\boldsymbol{X}_{1})\right),\\
\sigma_{33} & = & \sigma_{F}^{2} & = & 4\Var\mu_{wf}(\boldsymbol{X}_{1}) & ,\quad & \sigma_{23} & = & \sigma_{CF}^{2} & = & 4\Cov\left(\mu_{wc}(\boldsymbol{X}_{1}),\mu_{wf}(\boldsymbol{X}_{1})\right).
\end{alignat*}
Then 
\[
\sqrt{n}(\hat{\boldsymbol{p}}-\boldsymbol{p})\stackrel{d}{\to}N(0,\Sigma).
\]
\end{lem}

\begin{proof}
See Lemma 1 of \citet{Anon2022-ep}. The definitions of $\mu_{wa}(\boldsymbol{X}_{1}),\mu_{wc}(\boldsymbol{X}_{1})$
and $\mu_{wf}(\boldsymbol{X}_{1})$ can be found there.
\end{proof}
An application of the delta method yields the following.
\begin{prop}
\label{prop:The-asymptotic-variances}The coefficients in Table \ref{tab:coefficients}
are asymptotically normal, and their asymptotic variances are
\begin{eqnarray}
\sigma_{F}^{2} & = & \sigma_{A}^{2}\frac{1}{(1-p_{wf})^{2}}-2\sigma_{FA}\frac{1-p_{wa}}{(1-p_{wf})^{3}}+\sigma_{F}^{2}\frac{(1-p_{wa})^{2}}{(1-p_{wf})^{4}}.\label{eq:fleiss variance}\\
\sigma_{C}^{2} & = & \sigma_{A}^{2}\frac{1}{(1-p_{wc})^{2}}-2\sigma_{CA}\frac{1-p_{wa}}{(1-p_{wc})^{3}}+\sigma_{C}^{2}\frac{(1-p_{wa})^{2}}{(1-p_{wc})^{4}},\label{eq:cohen variance}\\
\sigma_{BP}^{2} & = & \sigma_{A}^{2}\frac{C^{2}}{(1-C)^{2}},\label{eq:brennan-prediger variance}\\
\sigma_{CF}^{2} & = & (1-p_{wf})^{-2}\left(1,-1,\nu_{CBP}\right)\Sigma\left(1,-1,\nu_{CBP}\right)^{T},\label{eq:cohen-fleiss variance}\\
\sigma_{CBP}^{2} & = & \frac{\sigma_{A}^{2}-2\sigma_{CA}+\sigma_{C}^{2}}{(1-\boldsymbol{1}^{T}W\boldsymbol{1}/C^{2})^{2}}.\label{eq:cohen-brennan-prediger variance}
\end{eqnarray}
\end{prop}

\begin{proof}
The expressions for $\sigma_{F}^{2}$ and $\sigma_{C}^{2}$ are from
\citet[Proposition 2]{Anon2022-ep}. The simple proof for the Cohen--Fleiss
coefficient is in the appendix, p. \ref{subsec:asymptotic variances}.
The asymptotic variance of the Brennan--Prediger coefficient is well-known
and the easiest to derive. The variance of the Cohen--Brennan--Prediger
coefficient follows immediately from an application of the delta method.
\end{proof}
To estimate the $\sigma_{x}^{2}$, where $x$ is a placeholder for
$F,C,BP,CF,$or $CBP$, we use an empirical approach that coincides
with that of \citet{Gwet2021-ht} in the special case of Fleiss' kappa
with nominal weights. See the comments following Proposition 1 of
\citet{Anon2022-ep} for details.

\subsection{Confidence intervals}

\citet{Anon2022-ep} found that the arcsine interval tends to do slightly
better than the untransformed interval for agreement coefficients
with rectangular design. For that reason, we will only look at the
arcsine interval here. Using the delta method, together with the fact
that $\frac{d}{dx}\arcsin(x)=1/\sqrt{1-x^{2}}$, we find that
\begin{equation}
\sqrt{n}(\arcsin\hat{\nu}_{x}-\arcsin\nu_{x})\stackrel{d}{\to}N(0,(1-\nu_{x}^{2})^{-1}\sigma_{x}^{2}),\label{eq:arcsine interval}
\end{equation}
where $\stackrel{d}{\to}$ denotes convergence in distribution and
$\arcsin x$ is the inverse of the sine function. Again, $x$ is a
placeholder for $F,C,BP,CF,$or $CBP$. Define the arcsine interval
as 
\begin{equation}
CI_{x}=\sin\left(\arcsin\hat{\nu}\pm t_{1-\alpha/2}(n-1)(1-\hat{\nu}_{x}^{2})^{-1}\hat{\sigma}_{x}^{2}\right),\label{eq:arcsine CI}
\end{equation}
where $\hat{\sigma}_{x}^{2}$ is estimated empirically, as described
in the previous subsection, and $\hat{\nu}_{x}$ is the sample estimator
of $\nu_{x}$.
\begin{example}[Example \ref{exa:Zapf exampe} (cont.)]
\label{exa:Zapf 2016 (cont)} We calculate arcsine confidence intervals
along with the point estimates for the five coefficients using the
data from \citet{Zapf2016-eu}. The results are in Table \ref{tab:confidence-zapf}.
\begin{table}
\noindent \centering{}\caption{\label{tab:confidence-zapf}Confidence limits for \citet{Zapf2016-eu}}
\begin{tabular}{cccccc}
\toprule & $\nu_{CF}$ & $\nu_{F}$ & $\nu_{C}$ & $\nu_{BP}$ & $\nu_{CBP}$\tabularnewline
\midrule Upper limit & $0.68$ & $0.67$ & $0.67$ & $0.70$ & $0.62$\tabularnewline
Lower limit & $0.46$ & $0.44$ & $0.45$ & $0.49$ & $0.41$\tabularnewline
Estimate & $0.57$ & $0.56$ & $0.57$ & $0.60$ & $0.52$\tabularnewline
\bottomrule &  &  &  &  & \tabularnewline
\end{tabular}
\end{table}
\end{example}


\subsection{Coverage of the confidence intervals}

We use the arcsine interval and nominally weighted coefficients. The
settings of our simulation study follows the settings of the first
sensitivity analysis closely. We use $N=10,000$ repetitions and the
following simulation parameters:
\begin{enumerate}
\item \textbf{Number of judges} $R$. We use $2,5,20$, corresponding to
a small, medium, and large selection of judges.
\item \textbf{Sample sizes} $n$. We use $n=20,100$, corresponding to small
and large agreement studies. 
\item \textbf{Model.} We simulate from the judge skill model used in the
sensitivity study (p. \pageref{subsec:sensitivity independence}).
\end{enumerate}
In some cases the simulation yields data frames with only identical
values. Our confidence interval construction do not cover these instances,
so we decided to discard these simulations, repeating the simulation
until we got a data frame with at least two different values.

\subsubsection{True distribution centered on the uniform distribution}

We use the same setup as in \ref{subsec:True-distribution-centered-uniform},
where we studied deviations from two assumptions. (a), that the true
distribution is centered on the uniform distribution, (b) that the
guessing distributions are equal. 

\begin{table}
\caption{\label{tab:coverage-marginal}Coverage and lengths of confidence intervals,
deviation from uniform.}

\noindent \centering{}%
\begin{tabular}{cccccccccccc}
\toprule & Coefficient & $\Var t$: & \multicolumn{3}{c}{None} & \multicolumn{3}{c}{Low} & \multicolumn{3}{c}{High}\tabularnewline
\cmidrule{3-12} &  & $\Var q$: & None & Low & High & None & Low & High & None & Low & High\tabularnewline
\midrule  &  & \multicolumn{10}{c}{$n=20$}\tabularnewline
 & Cohen--Fleiss &  & 0.95 & 0.95 & 0.95 & 0.95 & 0.95 & 0.95 & 0.82 & 0.81 & 0.82\tabularnewline
 &  &  & 0.26 & 0.26 & 0.26 & 0.26 & 0.26 & 0.26 & 0.33 & 0.33 & 0.33\tabularnewline
 & Fleiss &  & 0.95 & 0.95 & 0.95 & 0.94 & 0.95 & 0.95 & 0.81 & 0.81 & 0.82\tabularnewline
 &  &  & 0.26 & 0.26 & 0.26 & 0.27 & 0.27 & 0.27 & 0.33 & 0.33 & 0.34\tabularnewline
 & Cohen &  & 0.95 & 0.95 & 0.95 & 0.95 & 0.95 & 0.95 & 0.81 & 0.81 & 0.82\tabularnewline
 &  &  & 0.26 & 0.26 & 0.26 & 0.27 & 0.27 & 0.26 & 0.33 & 0.33 & 0.33\tabularnewline
 & BP &  & 0.96 & 0.96 & 0.96 & 0.96 & 0.96 & 0.96 & 0.96 & 0.96 & 0.95\tabularnewline
 &  &  & 0.26 & 0.26 & 0.26 & 0.26 & 0.26 & 0.26 & 0.26 & 0.26 & 0.26\tabularnewline
 & Cohen--BP &  & 0.92 & 0.92 & 0.92 & 0.90 & 0.90 & 0.90 & 0.43 & 0.43 & 0.44\tabularnewline
 &  &  & 0.27 & 0.27 & 0.27 & 0.28 & 0.28 & 0.28 & 0.33 & 0.33 & 0.33\tabularnewline
\midrule  &  & \multicolumn{10}{c}{$n=100$}\tabularnewline
 & Cohen--Fleiss &  & 0.95 & 0.95 & 0.95 & 0.94 & 0.94 & 0.94 & 0.52 & 0.52 & 0.53\tabularnewline
 &  &  & 0.11 & 0.11 & 0.11 & 0.11 & 0.11 & 0.11 & 0.14 & 0.14 & 0.14\tabularnewline
 & Fleiss &  & 0.95 & 0.95 & 0.95 & 0.94 & 0.94 & 0.93 & 0.51 & 0.52 & 0.51\tabularnewline
 &  &  & 0.11 & 0.11 & 0.11 & 0.11 & 0.11 & 0.11 & 0.14 & 0.14 & 0.14\tabularnewline
 & Cohen &  & 0.95 & 0.95 & 0.95 & 0.94 & 0.94 & 0.94 & 0.51 & 0.52 & 0.52\tabularnewline
 &  &  & 0.11 & 0.11 & 0.11 & 0.11 & 0.11 & 0.11 & 0.14 & 0.14 & 0.14\tabularnewline
 & BP &  & 0.95 & 0.95 & 0.95 & 0.95 & 0.95 & 0.95 & 0.95 & 0.95 & 0.90\tabularnewline
 &  &  & 0.11 & 0.11 & 0.11 & 0.11 & 0.11 & 0.11 & 0.11 & 0.11 & 0.11\tabularnewline
 & Cohen--BP &  & 0.95 & 0.94 & 0.95 & 0.89 & 0.89 & 0.88 & 0.10 & 0.11 & 0.11\tabularnewline
 &  &  & 0.11 & 0.11 & 0.11 & 0.11 & 0.11 & 0.11 & 0.14 & 0.14 & 0.14\tabularnewline
\bottomrule &  &  &  &  &  &  &  &  &  &  & \tabularnewline
\end{tabular}
\end{table}

Table \ref{tab:coverage-marginal} contains the results of the simulation.
All coefficients, except the Brennan--Prediger coefficient, perform
poorly when $t$ is far from the uniform. The Cohen--Fleiss coefficient
has poor coverage when the true distribution is far away from the
marginal distribution; likewise for Fleiss' kappa and Cohen's kappa.
The Brennan--Prediger coefficient performs surprisingly well, with
far better coverage than Cohen's kappa, Fleiss' kappa, and the Cohen--Fleiss
coefficient. On the other hand, the Cohen--Brennan--Prediger coefficient
performs poorly. The coverage when $t$ is far from the uniform or
$n=100$ is unacceptably low for all coefficients except the Brennan--Prediger
coefficient. 

\subsubsection{True distribution centered on the marginal distribution}

We use the same setup as in \ref{subsec:True-distribution-centered-marginal},
where we studied deviations from the assumption that the true distribution
is centered on the marginal distribution and that the guessing distributions
are equal. 

\begin{table}
\caption{\label{tab:coverage-uniform}Coverage and lengths of confidence intervals,
deviation from uniform.}

\noindent \centering{}%
\begin{tabular}{cccccccccccc}
\toprule & Coefficient & $\Var t$: & \multicolumn{3}{c}{None} & \multicolumn{3}{c}{Low} & \multicolumn{3}{c}{High}\tabularnewline
\cmidrule{3-12} &  & $\Var q$: & None & Low & High & None & Low & High & None & Low & High\tabularnewline
\midrule  &  & \multicolumn{10}{c}{$n=20$}\tabularnewline
 & Cohen--Fleiss &  & 0.95 & 0.95 & 0.95 & 0.93 & 0.92 & 0.92 & 0.37 & 0.37 & 0.39\tabularnewline
 &  &  & 0.27 & 0.27 & 0.3 & 0.29 & 0.29 & 0.32 & 0.36 & 0.36 & 0.34\tabularnewline
 & Fleiss &  & 0.95 & 0.95 & 0.94 & 0.92 & 0.92 & 0.91 & 0.37 & 0.37 & 0.39\tabularnewline
 &  &  & 0.27 & 0.28 & 0.31 & 0.30 & 0.30 & 0.34 & 0.36 & 0.37 & 0.36\tabularnewline
 & Cohen &  & 0.95 & 0.95 & 0.94 & 0.92 & 0.92 & 0.91 & 0.37 & 0.37 & 0.38\tabularnewline
 &  &  & 0.27 & 0.27 & 0.30 & 0.29 & 0.30 & 0.33 & 0.36 & 0.36 & 0.34\tabularnewline
 & BP &  & 0.96 & 0.95 & 0.90 & 0.96 & 0.95 & 0.91 & 0.94 & 0.94 & 0.87\tabularnewline
 &  &  & 0.26 & 0.26 & 0.26 & 0.26 & 0.26 & 0.27 & 0.26 & 0.26 & 0.26\tabularnewline
 & Cohen--BP &  & 0.88 & 0.85 & 0.72 & 0.69 & 0.67 & 0.55 & 0.04 & 0.04 & 0.06\tabularnewline
 &  &  & 0.29 & 0.3 & 0.34 & 0.32 & 0.32 & 0.35 & 0.26 & 0.26 & 0.25\tabularnewline
\midrule  &  & \multicolumn{10}{c}{$n=100$}\tabularnewline
 & Cohen--Fleiss &  & 0.95 & 0.95 & 0.95 & 0.85 & 0.85 & 0.84 & 0.09 & 0.09 & 0.13\tabularnewline
 &  &  & 0.11 & 0.11 & 0.12 & 0.12 & 0.12 & 0.14 & 0.18 & 0.18 & 0.18\tabularnewline
 & Fleiss &  & 0.95 & 0.95 & 0.95 & 0.85 & 0.84 & 0.82 & 0.09 & 0.09 & 0.12\tabularnewline
 &  &  & 0.11 & 0.11 & 0.13 & 0.12 & 0.12 & 0.14 & 0.18 & 0.18 & 0.19\tabularnewline
 & Cohen &  & 0.95 & 0.95 & 0.95 & 0.85 & 0.85 & 0.83 & 0.09 & 0.09 & 0.12\tabularnewline
 &  &  & 0.11 & 0.11 & 0.13 & 0.12 & 0.12 & 0.14 & 0.18 & 0.18 & 0.18\tabularnewline
 & BP &  & 0.92 & 0.89 & 0.64 & 0.91 & 0.87 & 0.65 & 0.86 & 0.81 & 0.58\tabularnewline
 &  &  & 0.11 & 0.11 & 0.11 & 0.11 & 0.11 & 0.11 & 0.11 & 0.11 & 0.10\tabularnewline
 & Cohen--BP &  & 0.81 & 0.70 & 0.28 & 0.34 & 0.28 & 0.15 & 0.00 & 0.00 & 0.01\tabularnewline
 &  &  & 0.12 & 0.12 & 0.14 & 0.13 & 0.13 & 0.15 & 0.12 & 0.12 & 0.11\tabularnewline
\bottomrule &  &  &  &  &  &  &  &  &  &  & \tabularnewline
\end{tabular}
\end{table}

The results are in Table \ref{tab:coverage-uniform}. The most striking
feature is, once again, the poor performance of every coefficient
except the Brennan--Prediger coefficient. The Brennan--Prediger
coefficient performs well, with a coverage of approximately $0.95$
in most scenarios when $n=20$. What's more, its length is always
the smallest. Cohen's kappa, Fleiss' kappa, and the Cohen--Fleiss
coefficient performs decently, except when the marginal distribution
is far away from the uniform, when their coverage is dismal. The Cohen--Brennan--Prediger
coefficients performs worse than the others, with larger confidence
interval lengths and horrible coverage. Compared to Table \ref{tab:coverage-marginal},
the coverages in Table \ref{tab:coverage-uniform} are much worse.
Even the best-performing Brennan--Prediger coefficient gets as low
as $0.58$ at a point. 

\section{Concluding remarks}

In the guessing model, a judge either knows -- with $100\%$ certainty
-- the correct classification, or he makes a guess. A more realistic
model would let knowledge come in degrees. A judge could be, say,
$70\%$ sure that a patient is X, $20\%$ sure that he is Y, and $10\%$
spread evenly on the remaining options. In epistemology, the \emph{credence
function} \citep{Pettigrew2019-li} quantifies his degree of belief
in the different propositions. Defining and working with knowledge
coefficients in more general ``credence models'' might be possible,
but identifiability issues looms large. 

The sensitivity and coverage studies in this paper are limited in
scope, as they only covers nominal weights and and a small number
of parameter tweaks. Larger simulation study could potentially confirm
or disconfirm our recommendation of reporting Cohen's kappa and Fleiss'
kappa together with the Brennan--Prediger coefficient. 

We reiterate that the agreement coefficient $\textrm{AC}_{1}$\citep{Gwet2008-qk}
is justified using a guessing model similar to the first Maxwell's
\citeyearpar{Maxwell1977-lf} (discussed on p. \pageref{subsec:maxwell}),
but it is not a knowledge coefficient. The relationship between the
$\textrm{AC}_{1}$ and the knowledge coefficient will be explored
in a future paper.

We have only discussed a rather peculiar sort of estimation of the
knowledge coefficient. It would, perhaps, be more natural to discuss
traditional estimation methods, such maximum likelihood estimation,
as explored by \citet{Aickin1990-nj} and \citet{Klauer1996-jd} in
their submodels of the guessing model. In particular, composite maximum
likelihood estimation \citep{Varin2011-nv} appears to be a good fit
to the problem. Bayesian estimation could also be a reasonable option.
If we only care about performance measures such as the mean squared
error, the Brennan--Prediger coefficient has small bias under many
scenarios, and its variance is virtually guaranteed to be smaller
than the variance of a composite maximum likelihood estimator. But
if we care about inference, even the superior confidence intervals
for the Brennan--Prediger coefficient have unacceptably poor coverage
under some circumstances. Since constructing approximate confidence
intervals for maximum likelihood and composite maximum likelihood
is routine, going this route will likely fix the coverage problem.

\subsection*{Declarations}

\textbf{Disclosure statement}. The authors do not have any conflicts
of interest to disclose. 

\bibliographystyle{apacite}
\bibliography{guessing}

\setcounter{secnumdepth}{0} 

\section*{Appendix\label{sec:Appendix}}

\subsection*{Proof of the Knowledge Coefficient Theorem on p. \pageref{thm:main theorem}.
\label{subsec:proof of knowledge theorem} }

Recall that $W$ is a matrix $C\times C$ matrix of agreement weights,
that is, a symmetric matrix whose elements $W_{ir}$ satisfy $1\geq W_{ir}$
with $W_{ir}=1$ if and only if $i=r$. We simplify our notation by
considering two judges only. Define $p_{wa}(r_{1},r_{2})$ and $p_{wc}(r_{1},r_{2})$
as the the probability of (chance) agreement restricted to the two
judges $r_{1}$ and $r_{2}$. Clearly, $p_{wa}=\binom{R}{2}^{-1}\sum_{r_{1}<r_{2}}p_{wa}(r_{1},r_{2})$
and $p_{wc}=\binom{R}{2}^{-1}\sum_{r_{1}<r_{2}}p_{wc}(r_{1},r_{2})$.

In the following lemma, we will view probability mass functions as
vectors. That is, we will view e.g. $p$ as the vector in $\mathbb{R}^{C}$
whose $i$th element is $p(x_{i})$. This will greatly simplify our
notation. 
\begin{lem}
\label{lem:expression lemma}The following is true.
\begin{enumerate}
\item The weighted agreement between $r_{1}$ and $r_{2}$, $p_{wa}(r_{1},r_{2})$,
equals 
\begin{multline}
E[s_{r_{1}}s_{r_{2}}]+(E[s_{r_{1}}]-E[s_{r_{1}}s_{r_{2}}])t^{T}Wq_{r_{2}}+(E[s_{r_{2}}]-E[s_{r_{1}}s_{r_{2}}])q_{r_{1}}^{T}Wt\\
+(1-E[s_{r_{1}}]-E[s_{r_{2}}]+E[s_{r_{1}}s_{r_{2}}])q_{r_{1}}^{T}Wq_{r_{2}}.\label{eq:p_wa_expression}
\end{multline}
\item The weighted chance agreement between $r_{1}$ and $r_{2}$, or $p_{wc}(r_{1},r_{2})$,
equals
\begin{multline}
E[s_{r_{1}}]E[s_{r_{2}}]t^{T}Wt+(E[s_{i_{1}r_{1}}]-E[s_{r_{1}}]E[s_{r_{2}}])t^{T}Wq_{r_{2}}+(E[s_{r_{2}}]-E[s_{r_{1}}]E[s_{r_{2}}])q_{r_{1}}^{T}Wt\\
+(1-E[s_{r_{1}}]-E[s_{r_{2}}]+E[s_{r_{1}}]E[s_{r_{2}}])q_{r_{1}}^{T}Wq_{r_{2}}.\label{eq:p^C_wa_expression}
\end{multline}
\item Finally, the weighted chance agreement between $r_{1}$ and $r_{2}$
can be written in terms of the marginal distributions and the weighting
matrix,
\begin{equation}
p_{wc}(r_{1},r_{2})=\sum_{x_{1},x_{2}}w(x_{1},x_{2})p_{r_{1}}(x_{1})p_{r_{2}}(x_{2})=p_{r_{1}}^{T}Wp_{r_{2}}.\label{eq:eq:p^F_wa_expression}
\end{equation}
\end{enumerate}
\end{lem}

\begin{proof}
\textbf{(i)} We will make use the following expression for $p_{wa}(r_{1},r_{2})$:
\[
p_{wa}(r_{1},r_{2})=\sum_{x^{\star}}\sum_{x_{1}}\sum_{x_{2}}w(x_{1},x_{2})p(x_{r_{1}}\mid s,x^{\star})p(x_{r_{2}}\mid s,x^{\star})t(x^{\star}).
\]
Recall that $x^{\star}$ is the true rating, $x_{1}$ is the rating
by judge $1$, and $x_{2}$ the rating by judge $2$, and the expression
for the full guessing model, $p(x\mid s,x^{\star})=s_{r}1[x=x^{\star}]+(1-s_{r})q_{r}(x)$
of formula (\ref{eq:full guessing model}). We expand the right hand
side of the expression for $p_{wa}(r_{1},r_{2})$ above to obtain

\begin{align*}
p_{wa}(r_{1},r_{2}) & =\sum_{x^{\star}}\sum_{x_{1}}\sum_{x_{2}}w(x_{1},x_{2})s_{r_{1}}s_{r_{2}}t(x^{\star})1[x_{1}=x_{2}=x^{\star}] & (=A)\\
 & +\sum_{x^{\star}}\sum_{x_{1}}\sum_{x_{2}}w(x_{1},x_{2})s_{r_{1}}(1-s_{r_{2}})t(x^{\star})1[x_{1}=x^{\star}]q_{r_{2}}(x_{2}) & (=B)\\
 & +\sum_{x^{\star}}\sum_{x_{1}}\sum_{x_{2}}w(x_{1},x_{2})s_{r_{2}}(1-s_{r_{2}})q_{r_{1}}(x_{1})t(x^{\star})1[x_{2}=x^{\star}] & (=C)\\
 & +\sum_{x^{\star}}\sum_{x_{1}}\sum_{x_{2}}w(x_{1},x_{2})(1-s_{r_{1}})(1-s_{r_{2}})q_{r_{1}}(x_{1})q_{r_{2}}(x_{2})t(x^{\star}) & (=D)
\end{align*}

Now we sum over $x^{\star},x_{1},x_{2}$ for each of $(A)-(D)$. Starting
with $(A)$, recall that $w(x,x)=1$ for all $x$. Thus
\begin{eqnarray*}
A & = & \sum_{x_{1}}\sum_{x_{2}}w(x_{1},x_{2})[s_{r_{1}}s_{r_{2}}t(x^{\star})1[x_{1}=x_{2}=x^{\star}],\\
 & = & s_{r_{1}}s_{r_{2}}.
\end{eqnarray*}
Now consider $(B)$, where we must recall that $W$, the weighting
matrix, is the matrix with elements $W_{ir}=w(x_{i},x_{r}$). 
\begin{eqnarray*}
B & = & \sum_{x^{\star}}\sum_{x_{1}}\sum_{x_{2}}w(x_{1},x_{2})s_{r_{1}}(1-s_{r_{2}})t(x^{\star})1[x_{1}=x^{\star}]q_{r_{2}}(x_{2}),\\
 & = & \sum_{x_{1}}\sum_{x_{2}}w(x_{1},x_{2})s_{r_{1}}(1-s_{r_{2}})t(x_{1})q_{r_{2}}(x_{2}),\\
 & = & s_{r_{1}}(1-s_{r_{2}})t^{T}Wq_{r_{2}}.
\end{eqnarray*}
Likewise, $C=s_{r_{2}}(1-s_{r_{2}})q_{r_{1}}^{T}Wt$ and $D=(1-s_{r_{1}})(1-s_{r_{2}})q_{r_{1}}^{T}Wq_{r_{2}}$. 

Summing over $x_{1},x_{2}$, the expression becomes 
\begin{eqnarray*}
 & = & s_{r_{1}}s_{r_{2}}+s_{r_{1}}(1-s_{r_{2}})q_{r_{1}}^{T}Wt+s_{r_{2}}(1-s_{r_{2}})q_{r_{1}}(x_{1})t(x_{2})\\
 &  & +(1-s_{r_{1}})(1-s_{r_{2}})q_{r_{1}}^{T}Wq_{r_{2}},
\end{eqnarray*}
and \ref{eq:p_wa_expression} follows from taking expectations over
$s_{r_{1}},s_{r_{2}}$.

\textbf{(ii)} We proceed in the same way as we did in (i). 
\begin{eqnarray*}
p_{wc}(r_{1},r_{2}\mid s_{r_{1}},s_{r_{2}}) & = & \sum_{x_{1},x_{2}}w(x_{r_{1}},x_{r_{2}})p(x_{r_{1}}\mid s_{r_{1}})p(x_{r_{2}}\mid s_{r_{1}})\\
 & = & \sum_{x_{1}^{\star}}\sum_{x_{2}^{\star}}\sum_{x_{1},x_{2}}w(x_{r_{1}},x_{r_{2}})p(x_{r_{1}}\mid s_{r_{1}},x_{1}^{\star})p(x_{r_{2}}\mid s_{r_{2}},x_{2}^{\star})t(x_{1}^{\star})t(x_{2}^{\star})
\end{eqnarray*}
Again, we expand this expression to obtain

\begin{align*}
p_{wc}(r_{1},r_{2}\mid s_{r_{1}},s_{r_{2}}) & =\sum_{x_{1}^{\star}}\sum_{x_{2}^{\star}}\sum_{x_{1}}\sum_{x_{2}}w(x_{1},x_{2})s_{r_{1}}s'_{r_{2}}t(x_{1}^{\star})t(x_{2}^{\star})1[x_{1}=x_{1}^{\star}]1[x_{2}=x_{2}^{\star}] & (=A)\\
 & +\sum_{x_{1}^{\star}}\sum_{x_{2}^{\star}}\sum_{x_{1}}\sum_{x_{2}}w(x_{1},x_{2})s_{r_{1}}(1-s{}_{r_{2}})t(x_{1}^{\star})1[x_{1}=x^{\star}]q_{r_{2}}(x_{2})t(x_{2}^{\star}) & (=B)\\
 & +\sum_{x_{1}^{\star}}\sum_{x_{2}^{\star}}\sum_{x_{1}}\sum_{x_{2}}w(x_{1},x_{2})s_{r_{2}}(1-s{}_{r_{1}})q_{r_{1}}(x_{1})t(x_{2}^{\star})1[x_{2}=x^{\star}]t(x_{1}^{\star}) & (=C)\\
 & +\sum_{x_{1}^{\star}}\sum_{x_{2}^{\star}}\sum_{x_{1}}\sum_{x_{2}}w(x_{1},x_{2})(1-s_{r_{1}})(1-s{}_{r_{2}})q_{r_{1}}(x_{1})q_{r_{2}}(x_{2})t(x_{1}^{\star})t(x_{2}^{\star}) & (=D)
\end{align*}
The main difference from (i) is in (A), 
\begin{eqnarray*}
 &  & \sum_{x_{1}^{\star}}\sum_{x_{2}^{\star}}\sum_{x_{1}}\sum_{x_{2}}w(x_{1},x_{2})s_{r_{1}}s{}_{r_{2}}t(x_{1}^{\star})t(x_{2}^{\star})1[x_{1}=x_{1}^{\star}]1[x_{2}=x_{2}^{\star}]\\
 & = & s_{r_{1}}s'_{r_{2}}t^{T}Wt.
\end{eqnarray*}
Let's consider (B) too:
\begin{eqnarray*}
 &  & \sum_{x_{1}^{\star}}\sum_{x_{2}^{\star}}\sum_{x_{1}}\sum_{x_{2}}w(x_{1},x_{2})s_{r_{1}}(1-s{}_{r_{2}})t(x_{1}^{\star})1[x_{1}=x^{\star}]q_{r_{2}}(x_{2})t(x_{2}^{\star}),\\
 & = & \sum_{x_{1}}\sum_{x_{2}}w(x_{1},x_{2})s_{r_{1}}(1-s{}_{r_{2}})t(x_{1})q_{r_{2}}(x_{2}),\\
 & = & s_{r_{1}}(1-s_{r_{2}})t^{T}Wq_{r_{2}}.
\end{eqnarray*}
(C) and (D) can be calculated in the same way. After taking expectations
with respect to independent $s_{r_{1}},s_{r_{2}}$, we find the expression
for $p_{wc}(r_{1},r_{2})$ in the statement of the Lemma.

\textbf{(iii)} The expression for $p_{wc}$ in terms of $p_{r_{1}},p_{r_{2}}$is
trivial.
\end{proof}
%
\begin{proof}[Proof of Theorem \ref{thm:main theorem}]

\textbf{(i).} Assume all guessing distributions are equal, i.e., $q_{r}(x)=q(x)$.
We wish to show that the following are equivalent
\begin{equation}
q(x)=t(x),\quad p(x)=t(x),\quad q(x)=p(x).\label{eq:equivalence equaitons}
\end{equation}
To do this, recall the marginal univariate model $p(x\mid s)=s_{r}t(x)+(1-s_{r})q(x).$
Take expectations over $s$ to obtain, with $\alpha=R^{-1}\sum_{r}E(s_{r}),$
\begin{equation}
p(x)=\alpha t(x)+(1-\alpha)q(x).\label{eq:marginal univariate model}
\end{equation}
It immediately follows that the expressions in \ref{eq:equivalence equaitons}
are equivalent.

Let's proceed to prove the rest of (i). If all guessing distributions
are equal to the true distribution, then the formula \ref{eq:p_wa_expression}
can be written as
\begin{multline*}
p_{wa}(r_{1},r_{2})=E[s_{r_{1}}s_{r_{2}}]+(E[s_{r_{1}}]-E[s_{r_{1}}s_{r_{2}}])t^{T}Wt+(E[s_{r_{2}}]-E[s_{r_{1}}s_{r_{2}}])t^{T}Wt\\
+(1-E[s_{r_{1}}]-E[s_{r_{2}}]+E[s_{r_{1}}s_{r_{2}}])t^{T}Wt.
\end{multline*}
Some of the terms cancel, leaving

\[
p_{wa}(r_{1},r_{2})=E[s_{r_{1}}s_{r_{2}}](1-t^{T}Wt)+t^{T}Wt.
\]
Moreover, the formula for $p_{wc}(r_{1},r_{2})$ can be simplified:
\begin{multline*}
p_{wc}(r_{1},r_{2})=E[s_{r_{1}}]E[s_{r_{2}}]t^{T}Wt+(E[s_{r_{1}}]-E[s_{r_{1}}]E[s_{r_{2}}])t^{T}Wt+(E[s_{r_{2}}]-E[s_{r_{1}}]E[s_{r_{2}}])t^{T}Wt\\
+(1-E[s_{r_{1}}]-E[s_{r_{2}}]+E[s_{r_{1}}]E[s_{r_{2}}])t^{T}Wt.
\end{multline*}
Most of the terms cancel, leaving only 
\[
p_{wc}(r_{1},r_{2})=t^{T}Wt.
\]
To verify these formulas, simply replace all instances of $q_{r}$
with $t$ in Lemma \ref{lem:expression lemma}, parts (i) and (ii).
Since the marginal distribution for every judge is the same under
the assumption that $q(x)=t(x)$, it follows that $p_{wa}=p_{wa}$.
Since the true distribution equals the marginal distribution,
\[
t^{T}Wt=p^{T}Wp=(R^{-1}\sum_{r}p_{r})^{T}W(R^{-1}\sum_{r}p_{r})=R^{-2}\sum_{r_{1},r_{2}}p_{r_{1}}^{T}Wp_{r_{2}},
\]
and by part (iii) of Lemma \ref{lem:expression lemma}, $p_{wf}=R^{-2}\sum_{r_{1},r_{2}}p_{r_{1}}^{T}Wp_{r_{2}}$,
Take the mean over all combinations of judges and reorder to arrive
at
\[
\nu=\frac{p_{wa}-p_{wc}}{1-t^{T}Wt}=\frac{p_{wa}-p_{wc}}{1-p_{wc}}=\frac{p_{wa}-p_{wf}}{1-p_{wf}}=\frac{p_{wa}-p_{wf}}{1-p_{wc}}.
\]

\textbf{(ii).} Assume that all guessing distributions are uniform.
Then if $q_{r}=C^{-1}\mathbf{1}$, which implies that $q_{r}^{T}p=C^{-1}$
for any probability mass function $p$. (This happens since $p$ sums
to $1$.) It follows that $q_{r_{2}}^{T}t=q_{r_{1}}^{T}q_{r_{2}}=C^{-1}$
for all $r_{1},r_{2}$. Using expression (i) of the Lemma and $W=I$,
we find that
\begin{multline*}
p_{a}(r_{1},r_{2})=E[s_{r_{1}}s_{r_{2}}]+(E[s_{r_{1}}]-E[s_{r_{1}}s_{r_{2}}])C^{-1}+(E[s_{r_{2}}]-E[s_{r_{1}}s_{r_{2}}])C^{-1}\\
+(1-E[s_{r_{1}}]-E[s_{r_{2}}]+E[s_{r_{1}}s_{r_{2}}])C^{-1}.
\end{multline*}
Canceling terms, we find that $p_{a}(r_{1},r_{2})=E[s_{r_{1}}s_{r_{2}}](1-C^{-1})+C^{-1}$.
It follows that $\nu=\frac{p_{a}-C^{-1}}{1-C^{-1}}$.

\textbf{(iii).} Suppose that $E[s_{r_{1}}s_{r_{2}}]=E[s_{r_{1}}]E[s_{r_{2}}]$
for all $r_{1},r_{2}$. Now subtract $p_{wc}(r_{1},r_{2})$ from $p_{wa}(r_{1},r_{2})$,
using Lemma \ref{lem:expression lemma}, parts (i) and (ii). Most
of the terms cancel, leaving us with
\[
p_{wa}(r_{1},r_{2})-p_{wc}(r_{1},r_{2})=E[s_{r_{1}}s_{r_{2}}](1-t^{T}Wt).
\]
Take the mean over all combinations of judges and reorder to arrive
at
\[
\nu=\frac{p_{wa}-p_{wc}}{1-t^{T}Wt}.
\]
If the true distribution equals the marginal distribution, then $p_{wf}=p^{T}Wp=t^{T}Wt$,
as explained in (i), hence $\nu=\frac{p_{wa}-p_{wc}}{1-p_{wf}}$,
as claimed. On the other hand, if the true distribution is uniform,
$t=C^{-1}\mathbf{1}$, hence $\nu=\frac{p_{wa}-p_{wc}}{1-\mathbf{1}^{T}W\mathbf{1}/C^{-2}}$.
\end{proof}
%

\subsection*{Proof of the expression for $\sigma_{CF}$ in Proposition \ref{prop:The-asymptotic-variances}\label{subsec:asymptotic variances}}

First, let us recall the multidimensional delta method. Let $f:\mathbb{R}^{k}\to\mathbb{R}$
be continuously differentiable at $\theta$ and suppose that $\sqrt{n}(\hat{\theta}-\theta)\stackrel{d}{\to}N(0,\Sigma)$.
Then 
\begin{equation}
\sqrt{n}[f(\hat{\theta})-f(\theta)]\stackrel{d}{\to}N(0,\nabla f(\theta)^{T}\Sigma\nabla f(\theta))\label{eq:delta method}
\end{equation}
In the case of the Cohen--Fleiss coefficient, $\theta=(p_{wa},p_{wc},p_{wf})$
and $f(\theta)=\frac{p_{wa}-p_{wc}}{1-p_{wf}}$. Then

\[
\nabla f=\frac{1}{1-p_{wf}}\left(1,-1,\nu_{CBP}\right).
\]

\[
(p_{wa},p_{wc},p_{wf})
\]
Thus the variance is 
\begin{eqnarray*}
(1-p_{wf})^{2}\nabla f(\theta)^{T}\Sigma\nabla f(\theta) & = & \left(1,-1,\nu_{CBP}\right)^{T}\Sigma\left(1,-1,\nu_{CBP}\right),\\
 & = & p_{wa}\sigma_{wa}^{2}+p_{wc}^{2}\sigma_{wc}^{2}+p_{wf}^{2}\sigma_{wf}^{2}+2\sigma_{wf}^{2}\sigma_{wc}^{2}.
\end{eqnarray*}


\subsection*{Proof that Aickin's model is a guessing model\label{subsec:Proof-that-Aickin's}}
\begin{lem}
Let the number of judges be $2$ and $s\sim\text{Bernoulli}(\alpha)$
be the same for both judges. Then the guessing model has unconditional
distribution
\begin{equation}
p(x_{1},x_{2})=\alpha1[x_{1}=x_{2}]t(x_{1})+(1-\alpha)q_{1}(x_{1})q_{2}(x_{2}).\label{eq:generalized aicken model}
\end{equation}
\end{lem}

\begin{proof}
Expanding the guessing model \ref{eq:full guessing model}
\begin{eqnarray*}
p(x_{1},x_{2}\mid s,x^{\star}) & = & s^{2}1[x_{1}=x^{\star}]1[x_{2}=x^{\star}]\\
 & = & +s(1-s)1[x_{1}=x^{\star}]q_{2}(x_{2})\\
 &  & +s(1-s)1[x_{2}=x^{\star}]q_{1}(x_{1})\\
 &  & +(1-s)(1-s)q_{1}(x_{1})q_{2}(x_{2})
\end{eqnarray*}
Since $s\sim\text{Bernoulli}(\alpha)$, we have that $Es^{2}=\alpha$,
$E(s(1-s))=0$ and $E(1-s)(1-s)=1-\alpha$. It follows that
\[
p(x_{1},x_{2}\mid x^{\star})=\alpha1[x_{1}=x^{\star}]1[x_{2}=x^{\star}]+(1-\alpha)q_{1}(x_{1})q_{2}(x_{2}).
\]
Summing over $x^{\star}$ yields $p(x_{1},x_{2})=\alpha1[x_{1}=x_{2}]t(x_{1})+(1-\alpha)q_{1}(x_{1})q_{2}(x_{2})$.
\end{proof}

\section*{Online appendix\label{sec:Online-appendix}}

\subsection*{More sensitivity analysis \label{subsec:more senstivity}}

This is a continuation of the sensitivity analysis of section \pageref{subsec:dependence}
when $\rho=0.5$ and $\rho=0.9$. If you compare these Tables with
Tables \ref{tab:centered on uniform0.2} (p. \vpageref{tab:centered on uniform0.2})
and \ref{tab:centered on marginal0.2} (p. \vpageref{tab:centered on marginal0.2})
you will find few differences. 

\begin{table}
\caption{\label{tab:centered on uniform-1-1}Sensitivity analysis. True distribution
centered on the uniform distribution ($\rho=0.9$).}

\noindent \begin{centering}
\begin{tabular}{ccccccccc}
\toprule & \multicolumn{2}{c}{Variability} &  & \multicolumn{5}{c}{Coefficient}\tabularnewline
\cmidrule[0.5pt]{2-3}\cmidrule[0.5pt]{5-9} & Guessing$^{\star}$ & True$^{*}$ &  & Cohen-Fleiss & Fleiss & Cohen & BP & Cohen--BP\tabularnewline
\midrule & \multirow{3}{*}{None} & None &  & \textbf{0.0e+00} & \textbf{0.0e+00} & \textbf{0.0e+00} & 1.1e-02 & 2.2e-02\tabularnewline
 &  & Low &  & \textbf{2.2e-02} & \textbf{2.2e-02} & \textbf{2.2e-02} & \textbf{1.2e-02} & \textbf{8.3e-02}\tabularnewline
 &  & High &  & 3.3e-01 & 3.3e-01 & 3.3e-01 & \textbf{1.8e-02} & 4.6e-01\tabularnewline
 & \multirow{3}{*}{Low} & None &  & \textbf{2.8e-05} & 3.4e-04 & 2.4e-04 & 1.4e-02 & 3.1e-02\tabularnewline
 &  & Low &  & \textbf{2.3e-02} & \textbf{2.3e-02} & \textbf{2.3e-02} & \textbf{1.4e-02} & \textbf{9.0e-02}\tabularnewline
 &  & High &  & 3.4e-01 & 3.4e-01 & 3.4e-01 & \textbf{2.2e-02} & 4.7e-01\tabularnewline
 & \multirow{3}{*}{High} & None &  & \textbf{2.3e-04} & 2.8e-03 & 1.9e-03 & 3.1e-02 & 7.1e-02\tabularnewline
 &  & Low &  & \textbf{2.4e-02} & \textbf{2.7e-02} & \textbf{2.6e-02} & \textbf{3.1e-02} & 1.3e-01\tabularnewline
 &  & High &  & 3.4e-01 & 3.5e-01 & 3.4e-01 & \textbf{3.7e-02} & 4.7e-01\tabularnewline
\bottomrule &  &  &  &  &  &  &  & \tabularnewline
\end{tabular}
\par\end{centering}
\vspace{1mm}

{\scriptsize{}$^{*}$Variability of the true distributions: Baseline:
True distribution is uniform.}{\scriptsize\par}

{\scriptsize{}$^{\star}$Variability of the guessing distributions.
Baseline: All guessing distributions are equal to the true distribution.}{\scriptsize\par}
\end{table}

\begin{table}
\caption{\label{tab:centered on marginal-1-1}Sensitivity analysis. True distribution
centered on the marginal guessing distribution ($\rho=0.9$).}

\noindent \begin{centering}
\begin{tabular}{ccccccccc}
\toprule & \multicolumn{2}{c}{Variability} &  & \multicolumn{5}{c}{Coefficient}\tabularnewline
\cmidrule[0.5pt]{2-3}\cmidrule[0.5pt]{5-9} & Guessing$^{\star}$ & True$^{*}$ &  & Cohen--Fleiss & Fleiss & Cohen & BP & Cohen--BP\tabularnewline
\midrule & \multirow{3}{*}{None} & None &  & \textbf{0.0e+00} & \textbf{0.0e+00} & \textbf{0.0e+00} & \textbf{0.0e+00} & \textbf{0.0e+00}\tabularnewline
 &  & Low &  & 3.9e-03 & 4.0e-03 & 3.9e-03 & \textbf{0.0e+00} & 1.2e-02\tabularnewline
 &  & High &  & 7.7e-02 & 7.8e-02 & 7.7e-02 & \textbf{0.0e+00} & 1.7e-01\tabularnewline
 & \multirow{3}{*}{Low} & None &  & 5.5e-05 & 4.4e-05 & 2.8e-05 & 4.4e-05 & \textbf{0.0e+00}\tabularnewline
 &  & Low &  & 3.7e-03 & 3.8e-03 & 3.8e-03 & \textbf{7.3e-04} & 1.2e-02\tabularnewline
 &  & High &  & 7.5e-02 & 7.6e-02 & 7.5e-02 & \textbf{2.5e-03} & 1.7e-01\tabularnewline
 & \multirow{3}{*}{High} & None &  & 8.3e-04 & 5.9e-04 & 4.0e-04 & 6.3e-04 & \textbf{0.0e+00}\tabularnewline
 &  & Low &  & \textbf{3.6e-03} & \textbf{4.5e-03} & \textbf{4.1e-03} & \textbf{2.8e-03} & 1.2e-02\tabularnewline
 &  & High &  & \textbf{7.4e-02} & \textbf{7.6e-02} & \textbf{7.5e-02} & \textbf{1.1e-02} & 1.7e-01\tabularnewline
\bottomrule &  &  &  &  &  &  &  & \tabularnewline
\end{tabular}
\par\end{centering}
\vspace{1mm}

{\scriptsize{}$^{*}$Variability of the true distributions: Baseline:
True distribution equals the marginal guessing distribution.}{\scriptsize\par}

{\scriptsize{}$^{\star}$Variability of the guessing distributions.
Baseline: All guessing distributions are equal.}{\scriptsize\par}
\end{table}

\begin{table}
\caption{\label{tab:centered on uniform-1-1-1}Sensitivity analysis. True distribution
centered on the uniform distribution ($\rho=0.5$).}

\noindent \begin{centering}
\begin{tabular}{ccccccccc}
\toprule & \multicolumn{2}{c}{Variability} &  & \multicolumn{5}{c}{Coefficient}\tabularnewline
\cmidrule[0.5pt]{2-3}\cmidrule[0.5pt]{5-9} & Guessing$^{\star}$ & True$^{*}$ &  & Cohen--Fleiss & Fleiss & Cohen & BP & Cohen--BP\tabularnewline
\midrule & \multirow{3}{*}{None} & None &  & \textbf{0.0e+00} & \textbf{0.0e+00} & \textbf{0.0e+00} & 1.1e-02 & 2.3e-02\tabularnewline
 &  & Low &  & \textbf{2.2e-02} & \textbf{2.2e-02} & \textbf{2.2e-02} & \textbf{1.2e-02} & \textbf{8.2e-02}\tabularnewline
 &  & High &  & 3.4e-01 & 3.4e-01 & 3.4e-01 & \textbf{1.9e-02} & 4.6e-01\tabularnewline
 & \multirow{3}{*}{Low} & None &  & \textbf{0.0e+00} & 3.9e-04 & 2.5e-04 & 1.5e-02 & 3.2e-02\tabularnewline
 &  & Low &  & \textbf{2.3e-02} & \textbf{2.3e-02} & \textbf{2.3e-02} & \textbf{1.6e-02} & \textbf{9.2e-02}\tabularnewline
 &  & High &  & 3.3e-01 & 3.4e-01 & 3.4e-01 & \textbf{2.3e-02} & 4.6e-01\tabularnewline
 & \multirow{3}{*}{High} & None &  & \textbf{0.0e+00} & 3.0e-03 & 1.9e-03 & 3.8e-02 & 8.8e-02\tabularnewline
 &  & Low &  & \textbf{2.4e-02} & \textbf{2.7e-02} & \textbf{2.6e-02} & \textbf{4.0e-02} & 1.4e-01\tabularnewline
 &  & High &  & 3.4e-01 & 3.5e-01 & 3.4e-01 & \textbf{4.7e-02} & 4.7e-01\tabularnewline
\bottomrule &  &  &  &  &  &  &  & \tabularnewline
\end{tabular}
\par\end{centering}
\vspace{1mm}

{\scriptsize{}$^{*}$Variability of the true distributions: Baseline:
True distribution is uniform.}{\scriptsize\par}

{\scriptsize{}$^{\star}$Variability of the guessing distributions.
Baseline: All guessing distributions are equal to the true distribution.}{\scriptsize\par}
\end{table}

\begin{table}
\caption{\label{tab:centered on marginal-1-1-1}Sensitivity analysis. True
distribution centered on the marginal guessing distribution ($\rho=0.5$).}

\noindent \begin{centering}
\begin{tabular}{ccccccccc}
\toprule & \multicolumn{2}{c}{Variability} &  & \multicolumn{5}{c}{Coefficient}\tabularnewline
\cmidrule[0.5pt]{2-3}\cmidrule[0.5pt]{5-9} & Guessing$^{\star}$ & True$^{*}$ &  & Cohen--Fleiss & Fleiss & Cohen & BP & Cohen--BP\tabularnewline
\midrule & \multirow{3}{*}{None} & None &  & \textbf{0.0e+00} & \textbf{0.0e+00} & \textbf{0.0e+00} & 6.60e-02 & 1.44e-01\tabularnewline
 &  & Low &  & 2.19e-01 & 2.19e-01 & 2.19e-01 & \textbf{4.89e-02} & 3.24e-01\tabularnewline
 &  & High &  & 5.95e-01 & 5.95e-01 & 5.95e-01 & \textbf{4.88e-02} & 6.33e-01\tabularnewline
 & \multirow{3}{*}{Low} & None &  & \textbf{7.30e-05} & 3.81e-04 & 2.84e-04 & 4.90e-02 & 1.08e-01\tabularnewline
 &  & Low &  & \textbf{2.99e-02} & \textbf{3.02e-02} & \textbf{3.01e-02} & \textbf{5.05e-02} & 1.65e-01\tabularnewline
 &  & High &  & 3.47e-01 & 3.47e-01 & 3.47e-01 & \textbf{5.38e-02} & 4.83e-01\tabularnewline
 & \multirow{3}{*}{High} & None &  & \textbf{6.00e-04} & 3.13e-03 & 2.33e-03 & 6.48e-02 & 1.47e-01\tabularnewline
 &  & Low &  & \textbf{3.14e-02} & \textbf{3.36e-02} & \textbf{3.28e-02} & \textbf{6.47e-02} & 1.91e-01\tabularnewline
 &  & High &  & 3.47e-01 & 3.52e-01 & 3.48e-01 & \textbf{6.71e-02} & 4.93e-01\tabularnewline
\bottomrule &  &  &  &  &  &  &  & \tabularnewline
\end{tabular}
\par\end{centering}
\vspace{1mm}

{\scriptsize{}$^{*}$Variability of the true distributions: Baseline:
True distribution equals the marginal guessing distribution.}{\scriptsize\par}

{\scriptsize{}$^{\star}$Variability of the guessing distributions.
Baseline: All guessing distributions are equal.}{\scriptsize\par}
\end{table}


\subsection*{Exact variance of $\hat{p}_{a}$ and $\hat{\nu}_{BP}$ \label{subsec:variance of p_a}}

In this section we derive the exact, non-asymptotic variance of $\hat{p}_{a}$
and the Brennan--Prediger coefficient $\hat{\nu}_{BP}$.
\begin{prop}
\label{prop:exact variane}The variance of 
\[
\hat{p}_{a}=n^{-1}\sum_{i=1}^{n}\binom{R}{2}^{-1}\sum_{r>r'}1[X_{ir}=X_{ir'}]
\]
is $\Var\hat{p}_{a}=n^{-1}(p_{2a}-p_{a}^{2})$. Here 
\[
p_{2a}=\binom{R}{2}^{-2}\sum_{r_{1}>r_{1}'}\sum_{r_{2}>r_{2}'}\mu_{r_{1}r_{1}'r_{2}r_{2}'},
\]
and 
\[
\mu_{r_{1}r_{1}'r_{2}r_{2}'}=P(X_{ir_{1}}=X_{ir_{1}'},X_{ir_{2}}=X_{ir_{2}'}).
\]
\end{prop}

\begin{proof}
Define $\mu_{rr}'=P(X_{ir}=X_{ir'})$. Then 
\[
p_{a}=\binom{R}{2}^{-1}\sum_{r>r'}P(X_{ir}=X_{ir'})=\binom{R}{2}^{-1}\sum_{r>r'}\mu_{rr}'
\]
by definition and thus,
\[
p_{a}^{2}=\left[\binom{R}{2}^{-1}\sum_{r>r'}\mu_{rr}'\right]^{2}=\binom{R}{2}^{-2}\sum_{r_{1}>r_{1}'}\sum_{r_{2}>r_{2}'}\mu_{r_{1}r_{1}'}\mu_{r_{2}r_{2}'}.
\]
To find the variance we must calculate
\begin{eqnarray*}
E\hat{p}_{a}^{2} & = & n^{-2}\binom{R}{2}^{-2}\sum_{i_{1},i_{2}}\sum_{r_{1}>r_{1}'}\sum_{r_{2}>r_{2}'}E\left\{ 1[X_{i_{1}r_{1}}=X_{i_{1}r_{1}'}]1[X_{i_{2}r_{2}}=X_{i_{2}r_{2}'}]\right\} ,\\
 & = & n^{-2}\binom{R}{2}^{-2}\sum_{i_{1},i_{2}}\sum_{r_{1}>r_{1}'}\sum_{r_{2}>r_{2}'}P(X_{i_{1}r_{1}}=X_{i_{1}r_{1}'},X_{i_{2}r_{2}}=X_{i_{2}r_{2}'}).
\end{eqnarray*}
Looking at $P(X_{i_{1}r_{1}}=X_{i_{1}r_{1}'},X_{i_{2}r_{2}}=X_{i_{2}r_{2}'})$,
we notice that, if $i_{1}\neq i_{2}$ it equals $P(X_{i_{1}r_{1}}=X_{i_{1}r_{1}'},X_{i_{2}r_{2}}=X_{i_{2}r_{2}'})=\mu_{r_{1}r_{1}'}\mu_{r_{2}r_{2}'}$
due to independence. This happens $n(n-1)$ times when we sum over
$i_{1},i_{2}$. On the other hand, if $i_{1}=i_{2}$, it equals $\mu_{r_{1}r_{1}'r_{2}r_{2}'}$.
This happens $n$ times when we sum over $i_{1},i_{2}$.

It follows that 
\[
E\hat{p}_{a}^{2}=(1-n^{-1})p_{a}^{2}+n^{-1}\binom{R}{2}^{-2}\sum_{r_{1}>r_{1}'}\sum_{r_{2}>r_{2}'}\mu_{r_{1}r_{1}'r_{2}r_{2}'},
\]
and the variance equals
\[
\Var\hat{p}_{a}=n^{-1}(E\hat{p}_{a}^{2}-p_{a}^{2})=n^{-1}(p_{2a}-p_{a}^{2}),
\]
as claimed.
\end{proof}
The sample Brennan--Prediger coefficient $\hat{\nu}_{BP}=(\hat{p}_{a}-C^{-1})/(1-C^{-1})$
is easy to deal with, as it's a linear function of $\hat{p}_{a}$.
As shown in Proposition \pageref{prop:exact variane}, the sample
agreement $\hat{p}_{a}$ has exact variance $\Var(\hat{p}_{a})=n^{-1}(p_{2a}-p_{a}^{2})$,
where $p_{a}$ is the true agreement and$p_{2a}=\binom{R}{2}^{-2}\sum_{r_{1}>r_{1}'}\sum_{r_{2}>r_{2}'}P(X_{ir_{1}}=X_{ir_{1}'},X_{ir_{2}}=X_{ir_{2}'}).$
The value of $p_{2a}$ is easy to estimate using the sample estimator
$\hat{p}_{2a}$. Moreover, the expectation of $\hat{p}_{a}^{2}$ is
$\frac{n-1}{n}p_{a}^{2}+\frac{1}{n}p_{2a}$, hence a debiased estimator
of $n^{-1}(p_{2a}-p_{a}^{2})$ is
\[
\hat{\Var}\hat{p}_{a}=(n-1)^{-1}(\hat{p}_{2a}-\hat{p}_{a}^{2}),
\]
provided $\hat{p}_{2a}>\hat{p}_{a}^{2}$. This inequality is rarely
broken, even in small samples.

Using the expression for $\Var\hat{p}_{a}$, it follows that the exact
variance of the estimated Brennan--Prediger coefficient $\hat{\nu}$
is 
\begin{equation}
\Var\hat{\nu}=n^{-1}(p_{2a}-p_{a}^{2})\frac{C^{2}}{(C-1)^{2}},\label{eq:Brennan-Prediger variance}
\end{equation}
and $\sqrt{n}(\hat{\nu}-\nu)$ is asymptotically normal by the delta
method. A natural plug-in estimator of $\Var\hat{\nu}$ is then
\[
\hat{\Var}\hat{\nu}=(n-1)^{-1}(\hat{p}_{2a}-\hat{p}_{a}^{2})\frac{C^{2}}{(C-1)^{2}}
\]
The usual estimator of $\Var\hat{p}_{a}$ is $n^{-1}(n-1)^{-1}\sum(\hat{p}_{a}(i)-\hat{p}_{a})^{2}$,
where $\hat{p}_{a}(i)$ is the probability of agreement restricted
to the $i$th item. This estimator is consistent since $\hat{p}_{a}=\frac{1}{n}\sum_{i=1}^{n}\hat{p}_{a}(i)$.
We would recommend this estimator of the variance over the plug-in
estimator for three reasons. First, it performs better than the plugin
estimator in several scenarios. Second, it is always non-negative,
being zero if and only if all $\hat{p}_{ai}$s are equal. Finally,
it is considerably less computationally expensive. 
\end{document}
